\documentclass[11pt]{article}
\usepackage[a4paper,margin=1in]{geometry}
\usepackage{amsmath,amssymb,amsthm}
\usepackage{booktabs}
\usepackage{graphicx}
\usepackage{xcolor}
\usepackage{braket}
\usepackage{hyperref}
\hypersetup{colorlinks=true,linkcolor=blue,urlcolor=blue,citecolor=blue}
\newcommand{\rt}[1]{\texttt{#1}}
\setlength{\emergencystretch}{2em}
\newcommand{\fn}[1]{\texttt{#1}}
\newcommand{\ifgraphic}[2]{\IfFileExists{#1}{\includegraphics[width=#2]{#1}}%
  {\fbox{\parbox{#2}{\centering\small figure \texttt{\detokenize{#1}} pending}}}}

\title{\textbf{Report R32 --- Floquet spectra and thermalisation\\ of the sixteen unitary quantum cellular automata}\\
\large Krylov sector analysis of quantum cellular automata}
\author{Tier 1 analysis}
\date{2026-10-14, revised 2026-10-15}

\begin{document}
\maketitle

\begin{abstract}
We study whether the dynamics inside the Krylov sectors of the sixteen unitary controlled-Hadamard brickwork circuits
(the quantum elementary cellular automata of Hillberry \emph{et al.}) thermalise. The tools are symmetry-resolved
quasienergy statistics of the one-period operator $U$, eigenstate entanglement, the spectral form factor, ETH for a local
observable and quench dynamics. Open chains reach $N=14$ to $20$, depending on the rule; rings reach $N=16$.
\begin{itemize}
\item \emph{One number organises the rules:} the count $k$ of firing neighbour patterns.
\item \emph{$k=3$} (54, 57, 99, 147) is chaotic on both boundaries, with COE or CUE as the symmetry dictates.
\item \emph{The one-sided $k=2$ rules} are chaotic on rings, protected by a newly identified dressed-reflection time
  reversal, but form a non-ergodic tower of conserved boundary projectors on open chains.
\item \emph{XOR} (free fermions) and \emph{$k=1$} (exponential fragmentation) do not thermalise.
\item \emph{All exact scars} of the chaotic rules descend from one universal eigenstate, $|h_+\rangle^{\otimes N}$.
\item \emph{With the $X$ gate} the chaotic rule 54 is the integrable RCA54: quantum chaos follows the firing count, not
  the classical dynamics.
\end{itemize}
Sec.~\ref{sec:summary} summarises all results, together with what was known before and what is new.
\end{abstract}

\tableofcontents

% R32 opening chapter: detailed summary (2026-10-15).  Input by R32_floquet_spectra.tex.
% Every number here appears, with its provenance, in a later section or a generated table.
\section{Summary}
\label{sec:summary}

\subsection{Question and model}
The Krylov analysis of the earlier reports describes \emph{how} the Hilbert space of the sixteen unitary rules falls
apart. This report asks whether the dynamics \emph{inside} the pieces thermalise.

The circuits are the quantum elementary cellular automata of Hillberry and coworkers~\cite{hillberry2016,hillberry2021}
(Sec.~\ref{sec:lit-known}):
\begin{itemize}
\item On a chain of $N$ qubits, the Hadamard $H$ acts on site $j$ iff the neighbour pattern $(x_{j-1},x_{j+1})$ is one of
  the rule's \emph{firing patterns}.
\item Even sites are updated first, then odd sites (brickwork), and one period is $U=L_BL_A$.
\item Two boundaries are studied: open with ghost neighbours fixed to 0 (obc0), and rings of even length.
\item The $X$ gate turns every rule into a classical reversible automaton and serves as the comparison.
\end{itemize}

The tools are applied to $U$ itself, not to a Hamiltonian:
\begin{itemize}
\item the quasienergy gap ratio $\langle r\rangle$ after \emph{full} symmetry resolution (Krylov sectors, gate-algebra
  blocks, reflection or the half step $\Omega$ on the chain, and momentum with the half step $\Omega_T$ on the ring);
\item the eigenstate entanglement against the Page value;
\item the spectral form factor;
\item ETH for a local observable;
\item quench dynamics.
\end{itemize}

\subsection{The classification}
One combinatorial number organises the sixteen rules: the count $k$ of the four neighbour patterns that fire the gate,
and, for $k=2$, whether the firing set depends on one neighbour or on both (Table~\ref{tab:summary-classes}).
\begin{table}[h]
\centering\small
\begin{tabular}{p{0.13\textwidth}p{0.12\textwidth}p{0.31\textwidth}p{0.31\textwidth}}
\toprule
class & rules & open chain & ring \\
\midrule
$k=3$ & 54, 57, 99, 147 & \textbf{chaotic}: $\langle r\rangle=0.523$--$0.531$ at $N=14$ (GOE 0.531) & \textbf{chaotic}: 57/99 CUE
($0.5985$ at $N=16$), 147 COE ($0.528$), 54 COE but slow ($0.510$, rising) \\
$k=2$ one-sided & 60, 102, 153, 195 & tower of conserved boundary projectors; $\langle r\rangle\approx0.424$, flat for
$N=13\ldots16$ & \textbf{chaotic}, COE ($0.523$--$0.527$ at $N=16$), protected by a hidden time reversal \\
$k=2$ XOR & 105, 150 & arithmetic: exactly $3^m$ distinct levels for 150 ($N=6\ldots15$) & arithmetic (150 is free
fermions) \\
$k=1$ & 108, 156, 198, 201 & exponential fragmentation; 108/201 Poisson, 156/198 $\langle r\rangle\approx0.42$ (common
$N=19$) & unchanged \\
$k=0$, $4$ & 204, 51 & trivial ($U=1$, $U=H^{\otimes N}$) & trivial \\
\bottomrule
\end{tabular}
\caption{The sixteen rules with $V=H$. Each class is compared at one common $N$ (Sec.~\ref{sec:results}); ring values are
for the complex momentum sectors (Sec.~\ref{sec:pbc}).}
\label{tab:summary-classes}
\end{table}

\paragraph{Only the $k=3$ class is chaotic on both boundaries.} Its ensemble follows the symmetry. With open boundaries
every block carries the time reversal $KL_A$ (COE). On the ring, the mirror rules 54 and 147 keep a time reversal $KL_AR$
(COE). 57 and 99 have none inside a momentum sector, so they are CUE, with a GUE-shaped $P(r)$ (Sec.~\ref{sec:pbc}).
Without symmetry resolution, 147 looks integrable ($0.42$ against $0.53$).

\paragraph{The one-sided class is a boundary phenomenon.}
\begin{itemize}
\item \emph{Open chain:} the free end carries a chain of conserved projectors onto Hadamard eigenstates, which builds a
  tower of blocks of size $2^{N-1},2^{N-2},\ldots$. Inside the largest block the statistics plateau at
  $\langle r\rangle\approx0.424$ for $N=13\ldots16$: neither Poisson nor GOE, and not drifting.
\item \emph{Ring:} there is no tower and the rules are COE-chaotic, although they have no mirror symmetry. The time
  reversal is a newly identified \emph{dressed bond reflection}: the reflection is combined with a rotation of one
  sublattice into the Hadamard eigenbasis (Sec.~\ref{sec:answers}).
\item \emph{Dynamics:} random product states reach only 0.72 (60) and 0.91 (153) of the Page entanglement on the open
  chain, but 0.99 to 1.00 on the ring.
\end{itemize}

\paragraph{XOR and $k=1$ do not thermalise.} 150 with $H$ is a free-fermion circuit~\cite{hillberry2024}; its
arithmetic spectrum reproduces R8's level count exactly. The $k=1$ rules split into exponentially many small blocks.
\begin{itemize}
\item In 156/198 a block is a product of segments, each holding one mobile kink with the exact dispersion
  $\cos\theta=-\sin^2(\pi m/(L+1))$. The sums of these few levels give $\langle r\rangle\approx0.42$.
\item 108/201 have larger, many-body rooms and Poisson statistics.
\end{itemize}

\subsection{Thermalisation}
All diagnostics give the same ordering (Sec.~\ref{sec:thermal}, Table~\ref{tab:thermal}):
\begin{itemize}
\item \emph{57/99 thermalise fastest and most cleanly:} ETH fluctuations decay as $d^{-0.40}$, close to $d^{-1/2}$; the
  SFF approaches the ramp after $t_2=78$ periods at $N=16$; random states reach Page.
\item \emph{153, 147, 60 and 54} are chaotic on the ring but converge progressively more slowly. Their ETH exponents
  are $-0.24$ to $-0.13$, $t_2$ is 5 to 13 times longer, and the SFF has a large short-time excess.
\item \emph{XOR and $k=1$} stay far from thermal in every diagnostic.
\end{itemize}
The gap ratio is the most robust single diagnostic, but not a sufficient one: the chaotic rules contain non-thermal
eigenstates and slow states (next subsection).

\subsection{Exact structures behind the classes}
\begin{itemize}
\item \textbf{One universal scar.} $|h_+\rangle^{\otimes N}$, the product of $+1$ Hadamard eigenstates, is an eigenstate
  of \emph{every} controlled-Hadamard gate, for every rule and boundary. It explains all exact scars of the chaotic
  rules:
  \begin{itemize}
  \item the state itself ($\theta=0$);
  \item its projection off a frozen vacuum (the entangled singlet of 54, 60 and 102, with entanglement in closed form,
    decreasing with $N$);
  \item $|h_-\rangle$ on boundary sites whose gate is control-independent ($\theta=\pi$; the counts 2/2/4/1 follow
    exactly).
  \end{itemize}
  Close to this vacuum live the quasi-bound edge mode of 57 (quasienergy $\pm0.312\pi$) and the slowly relaxing defect
  states of 54 and 147.
\item \textbf{Flip partners.} $X^{\otimes N}U_W(H)X^{\otimes N}=U_{W'}(XHX)\cong U_{W'}(-H)$: the partner rule with a
  $-1$ on every firing. This explains why flip partners share their block structure but not their spectra, and the obc0
  boundary breaks the flip even for $X$.
\item \textbf{The $X$ gate.} Classically, 54/147 is the integrable RCA54, 201/108 is the Floquet-PXP automaton, and
  XOR and the one-sided rules are linear. X-102 is exactly the solvable deterministic Floquet quantum East
  model~\cite{bertini2024det}. Only 57/99 have orbits that fill a finite fraction of the space, and only on the open
  chain; on the ring a conserved domain-wall number confines them. \emph{Quantum chaos with $H$ therefore follows the
  firing count, not the classical dynamics}: the integrable RCA54 becomes a chaotic circuit with $H$.
\end{itemize}

\subsection{What is known and what is new}
\emph{Known:}
\begin{itemize}
\item the model and its numbering~\cite{hillberry2016,hillberry2021};
\item entanglement-spectrum (Brody) indications that $T_{13}$ and $T_{14}$ (our 57, 54) are chaotic and $T_1$ (201) is
  not~\cite{hillberry2021};
\item the free-fermion nature and the domain-wall charge of the Hadamard Goldilocks rule 150~\cite{hillberry2024};
\item the classical integrable parents, and the solvable DFQEM (= X-102)~\cite{buca2021,wilkinson2020,bertini2024det}.
\end{itemize}
\emph{New, as far as we found} (Sec.~\ref{sec:lit-known}):
\begin{itemize}
\item the symmetry-resolved quasienergy classification of all sixteen rules on both boundaries, organised by $k$;
\item the boundary origin of the one-sided class and its hidden time reversal;
\item the universal scar and its consequences;
\item the exact kink dispersion of 156/198;
\item the flip map;
\item the domain-wall charge of X-57;
\item the SFF, ETH and quench characterisation of all classes.
\end{itemize}

\subsection{Open questions}
\begin{enumerate}
\item Why the one-sided rules plateau at $\langle r\rangle\approx0.424$ on the open chain. The tower block is the same
  chain with one boundary site fixed in $|h_-\rangle$; some further structure seems to be at work.
\item The origin of the large short-time SFF excess and the slow ETH convergence of 54, 60 and 147, presumably tied to
  the $|h_+\rangle$ vacuum and its slow defects.
\item An analytic form of the 57 edge mode.
\item Whether 54 reaches COE on the ring (0.510 at $N=16$, still rising).
\end{enumerate}

\paragraph{Reading guide.}
\begin{itemize}
\item Sec.~\ref{sec:lit}: the tools for Floquet operators and the literature on these circuits.
\item Secs.~\ref{sec:methods}--\ref{sec:algorithm}: methods and the optimised pipeline.
\item Sec.~\ref{sec:results}: open-chain results.
\item Sec.~\ref{sec:pbc}: rings.
\item Sec.~\ref{sec:x}: the X gate.
\item Sec.~\ref{sec:thermal}: thermalisation.
\item Sec.~\ref{sec:answers}: the exact structures and the runs beyond the dense limit.
\item Sec.~\ref{sec:memory}: compute.
\end{itemize}


\section{Scope and history}
\label{sec:scope}
Task (2026-10-14):
\begin{enumerate}
\item review the literature on eigenstate entanglement, level statistics and scars for Floquet operators;
\item complete the spectral dataset for the 16 unitary rules;
\item per-rule figures;
\item interpretation;
\item a storyline draft and candidate thesis figures, to be shown before any hand-off;
\item the same for the X gate (optional).
\end{enumerate}
Requests added during the work:
\begin{itemize}
\item resolve all symmetries;
\item check the classes on rings;
\item compare each class at one common $N$;
\item resolve the open physics points and push past the dense-sector limit for the relevant cases, verifying every
  assumption;
\item add the spectral form factor, ETH and dynamics;
\item a second literature pass (Hillberry; the Floquet quantum East circuits).
\end{itemize}
Corrections made along the way are recorded where they occur (Secs.~\ref{sec:pbc}, \ref{sec:answers}) and in
\fn{NEXT\_STEPS.md}.

% R32 section 2: literature review (input by R32_floquet_spectra.tex)
\section{The tools, and what they mean for a Floquet operator}
\label{sec:lit}

Everything in this report is a property of one operator, the one-period
evolution $U=L_BL_A$ of the brickwork circuit, and not of a Hamiltonian. This
section collects what the literature says the three diagnostics mean in that
setting, and where the Hamiltonian intuition has to be dropped.

\subsection{Quasienergies, and why there is no ``middle of the spectrum''}

For a time-periodic Hamiltonian $H(t+T)=H(t)$ or a periodically applied circuit,
the stroboscopic dynamics is generated by the Floquet operator
$U=\mathcal{T}\exp[-i\int_0^T H(t)\,dt]$, here simply the product of the two
gate layers. Its eigenvalues lie on the unit circle,
$U\ket{\phi_n}=e^{-i\theta_n}\ket{\phi_n}$, and the quasienergies $\theta_n$ are
defined only modulo $2\pi$. There is no ordering, no ground state and no
temperature axis.

For a generic interacting driven system the long-time state is the
\emph{infinite-temperature} state. D'Alessio and Rigol showed by exact
diagonalisation and finite-size scaling that for any non-zero driving period the
Floquet operator of a generic lattice system acquires circular-ensemble
properties as the system grows. The Floquet Hamiltonian is then non-local and
the system heats to infinite temperature~\cite{dalessio2014}. Lazarides, Das and
Moessner reached the same conclusion from the structure of Floquet eigenstates,
which mix many eigenstates of the undriven Hamiltonian. They also showed that
integrable driven systems are instead described by a periodic generalised Gibbs
ensemble~\cite{lazarides2014}. Ponte, Chandran, Papi\'c and Abanin found that
ergodic systems delocalise in energy space and heat up under both local and
global driving~\cite{ponte2015annphys}. Disorder can prevent this: the Floquet
eigenstates of a many-body localised system obey an area law, and a delocalised
phase is separated from the MBL phase by a transition seen in the eigenstate
entanglement~\cite{ponte2015prl}.

\emph{Consequence for the figures.} The Floquet version of the eigenstate
thermalisation hypothesis
(ETH)~\cite{deutsch1991,srednicki1994,rigol2008,dalessio2016} says that every
eigenstate of a chaotic $U$ looks locally like the infinite-temperature state.
Plotted against $\theta$, the eigenstate entanglement of a thermalising circuit
is therefore \emph{flat}, a band at the random-state value at every
quasienergy. There is no Hamiltonian ``dome'' that peaks in the middle of the
spectrum. A non-thermal (scar-like) eigenstate shows up as an outlier
\emph{below} the band, at whatever quasienergy it happens to sit.

\subsection{Spectral statistics: circular ensembles and the gap ratio}

Dyson's circular ensembles are the random-matrix models for unitary
operators~\cite{dyson1962}: orthogonal (COE, $\beta=1$), unitary (CUE,
$\beta=2$) and symplectic (CSE, $\beta=4$). The Bohigas--Giannoni--Schmit
conjecture~\cite{bgs1984} assigns random-matrix level statistics to quantum
chaotic systems. The Berry--Tabor conjecture assigns Poisson statistics to
integrable ones, with uncorrelated, clustering levels~\cite{berrytabor1977}.
Which circular ensemble applies is fixed by the antiunitary symmetries \emph{of
$U$ itself}. This need not match the class of the instantaneous Hamiltonian:
Regnault and Nandkishore showed that a Floquet system can be in the unitary
class while its instantaneous Hamiltonian is always orthogonal, and vice versa.
They also showed that a drive can thermalise even when the instantaneous
Hamiltonian is integrable~\cite{regnault2016}.

The gap ratio $r_n=\min(s_n,s_{n+1})/\max(s_n,s_{n+1})$ of consecutive
spacings, introduced by Oganesyan and Huse~\cite{oganesyan2007}, is independent
of the local density of states and so needs no unfolding. Atas \emph{et al.}\
derived Wigner-like surmises for its distribution
\begin{equation}
  P_\beta(r)=\frac{2}{Z_\beta}\,\frac{(r+r^2)^\beta}{(1+r+r^2)^{1+3\beta/2}},
  \qquad r\in[0,1],
  \label{eq:surmise}
\end{equation}
with $Z_1=8/27$ and $Z_2=4\pi/(81\sqrt3)$. Their large-matrix values,
$\langle r\rangle=0.5307$ (GOE) and $0.5996$ (GUE), are the reference values
used here. Poisson statistics give $P_0(r)=2/(1+r)^2$ and
$\langle r\rangle=2\ln2-1=0.3863$~\cite{atas2013}. For a chaotic Floquet
operator the quasienergy density is uniform, so the ratio is doubly convenient.

Two caveats from the literature are central to this report:
\begin{itemize}
\item \textbf{Symmetry sectors must be resolved.} Superposing independent
  spectra from different symmetry sectors removes level repulsion between them
  and pulls $\langle r\rangle$ towards the Poisson value, even if every sector
  is chaotic. Giraud, Mac\'e, Vernier and Alet derived the gap-ratio surmises for
  such superpositions and used them to detect hidden symmetries, including in
  periodically driven spin chains~\cite{giraud2022}. ETH, too, is a statement
  per symmetry sector. Mondaini \emph{et al.}\ made exactly this point in their
  comment~\cite{mondaini2018} on a claimed ETH counterexample~\cite{shiraishi2017}.
\item \textbf{In fragmented systems the Krylov sectors play the role of
  symmetry sectors.} Hilbert-space fragmentation (HSF) splits the space into
  exponentially many invariant subspaces that are not labelled by conventional
  symmetries~\cite{sala2020,khemani2020}. In the commutant language of Moudgalya
  and Motrunich~\cite{moudgalya2022}, the invariant subspaces are the blocks of
  the algebra generated by the local terms (or gates). Its commutant has
  exponential dimension, and the product-basis Krylov sectors (``classical''
  fragmentation) can be refined further by quantum
  fragmentation~\cite{moudgalya2022}. Level statistics are meaningful only
  inside one such block. Pooling blocks measures their number, not their
  chaoticity.
\end{itemize}

\subsection{Eigenstate entanglement}

For a random pure state on $\mathbb{C}^m\otimes\mathbb{C}^n$ with $m\le n$,
Page's mean entanglement entropy is $S_\mathrm{Page}\approx\ln m-m/(2n)$, short
of the maximum by about half a unit of information~\cite{page1993}. Eigenstates
of chaotic Hamiltonians approach this value in the middle of the spectrum.
Conservation laws lower it: Vidmar and Rigol showed that for random states at
fixed particle number away from half filling, the average entropy falls short of
the maximum by an amount growing as the square root of the
volume~\cite{vidmar2017}. For a fragmented system the right benchmark is
therefore a random state \emph{inside the relevant block}, not $S_\mathrm{Page}$
of the full space.

Volume-law eigenstates signal thermalisation. Area-law or otherwise anomalously
weakly entangled eigenstates embedded in a thermal spectrum are the defining
signature of quantum many-body scars (QMBS). QMBS were identified as the origin
of the persistent revivals in the Rydberg-atom experiment~\cite{bernien2017},
via scarred eigenstates of the PXP model~\cite{turner2018}. Embedding
constructions produce them systematically~\cite{shiraishi2017}. Reviews cover
the semiclassical picture~\cite{serbyn2021}, exact
results~\cite{moudgalya2022review}, and the quasiparticle
perspective~\cite{chandran2023}. Scars also exist intrinsically in driven
systems: Sugiura, Kuwahara and Saito proved that a PXP-type Floquet model hosts
exact scar eigenstates of $U$ at specific driving periods, a rigorous breakdown
of Floquet ETH~\cite{sugiura2021}.

\emph{Degeneracies make single-eigenvector entanglement ill-defined.} If an
eigenvalue of $U$ is $g$-fold degenerate, any orthonormal basis of the
eigenspace is equally valid, and the entanglement of ``the'' eigenvectors
depends on that arbitrary choice. Earlier R8 found that 43--74\,\% of the stored
rule-150 eigenstates lie in such eigenspaces, with multiplicity up to 35. This
report therefore (i)~resolves all symmetries first, so the remaining basis
freedom is as small as possible, (ii)~records $g$ for every state, and (iii)~adds
a basis-independent quantity for every eigenspace: the Haar-averaged R\'enyi-2
entropy (Sec.~\ref{sec:methods}).

\subsection{Quantum circuits and automata in particular}

Brickwork Floquet circuits are the setting in which many-body quantum chaos has
been connected analytically to random-matrix theory. Kos, Ljubotina and Prosen
derived RMT spectral correlations for periodically kicked many-body
systems~\cite{kos2018}. Chan, De Luca and Chalker solved a Floquet circuit of
random two-site gates (depth two, like ours) in the large-$q$ limit, including
its spectral form factor and entanglement growth~\cite{chan2018}. The circuits
studied here are the opposite extreme: \emph{one fixed} controlled gate,
repeated.

Automaton-based circuits are a known source of non-ergodic eigenstates.
Gopalakrishnan and Zakirov constructed facilitated quantum cellular automata,
mostly from Clifford gates. Their eigenstate entanglement differs from generic
chaotic systems in all five models, and some models have exponentially many
eigenstates with disentangled sites, tied to reflection
symmetry~\cite{gopalakrishnan2018}. Iadecola and Vijay showed that deforming a
classical reversible cellular automaton gives generically non-integrable Floquet
systems that inherit exact non-thermal eigenstates from the automaton's periodic
orbits~\cite{iadecola2020}. The automata themselves (our $V=X$ circuits) include
the interacting integrable Rule~54~\cite{buca2021} and the Floquet-PXP Rule~201,
whose solitons are domain walls between three vacua~\cite{wilkinson2020}.

\subsection{Two structural facts used throughout}

\paragraph{$U$ is real orthogonal and a product of two reflections.} Both gates,
$H$ and $X$, are real symmetric involutions, so every controlled gate is one and
each layer $L_{A,B}$ is a real symmetric involution (a reflection). Hence
$U=L_BL_A$ is real orthogonal. Two consequences:
\begin{enumerate}
\item Complex conjugation $K$ commutes with $U$, so within every
  symmetry-resolved block the spectrum is symmetric under
  $\theta\to-\theta$. The independent spectral data are the $|\theta|\in[0,\pi]$.
  Ratios computed over the full circle would count each spacing twice and add
  spurious small gaps where mirror pairs approach $0$ and $\pi$.
\item By the two-subspace theorem~\cite{jordan1875,halmos1969}, the space
  decomposes into the intersections of the $\pm1$ eigenspaces of $L_A$ and
  $L_B$ (eigenvalues $\pm1$ of $U$) and two-dimensional subspaces on which $U$
  rotates by $\pm2\varphi_j$, where the $\varphi_j$ are the principal angles
  between the $+1$ eigenspaces of the two layers. For generic reflections these
  angles follow a real ($\beta=1$) Jacobi ensemble, whose bulk statistics are
  those of the GOE. We verified the reference value directly with the same
  routine used for the data: products of two Haar-random reflections of size
  $300$--$800$ give $\langle r\rangle=0.533\pm0.004$ on the half spectrum,
  compared with $0.535\pm0.005$ for a GOE control and $0.604\pm0.005$ for a GUE
  control.
\end{enumerate}
In addition $T=L_AK$ is an antiunitary with $TUT^{-1}=U^{-1}$ and $T^2=1$, a
time-reversal symmetry, and also points to the orthogonal class. The chaotic
reference is therefore $\langle r\rangle_\mathrm{GOE}=0.5307$ for every block that
$T$ maps to itself: all blocks with open boundaries, and the self-conjugate momentum
sectors on a ring. In a ring-momentum sector $j\ne-j$, $T$ maps the sector to its
partner $-j$, no antiunitary symmetry remains inside it, and the chaotic reference is
the CUE value $0.5996$.

\paragraph{The X gate makes $U$ a permutation.} With $V=X$ every basis state is
mapped to a single basis state. On a cycle of length $\ell$ the spectrum is the
$\ell$-th roots of unity, a picket fence with $r\equiv1$, and the eigenvectors
are Fourier modes along the cycle. Each Krylov sector of a permutation is a
single cycle, so this holds sector by sector. The commutant study R-T26 found
exactly this, $P(r)=\delta(r-1)$, for rule 156. Neither Poisson nor RMT applies,
and an X-gate $\langle r\rangle$ measures orbit arithmetic, not chaos.

% R32 section 3: methods (input by R32_floquet_spectra.tex)
\section{Methods}
\label{sec:methods}

\subsection{Model and conventions}
The sixteen unitary rules are the tuples $(r_{00},r_{01},r_{10},r_{11})\in\{I,V\}^4$, Wolfram numbers
51, 54, 57, 60, 99, 102, 105, 108, 147, 150, 153, 156, 195, 198, 201 and 204,
with $V=H$ throughout this report (Sec.~\ref{sec:x} for $V=X$). Boundary
conditions are open with fixed vacuum neighbours (obc0). One period is
$U=L_BL_A$, with $L_A$ the gates on the odd (1-based) sites, applied first, and
$L_B$ those on the even sites. The quasienergy convention is
$U\ket{\phi}=e^{-i\theta}\ket{\phi}$, $\theta\in(-\pi,\pi]$. The bipartition is
$A=$ sites $1\ldots\lfloor N/2\rfloor$.

\subsection{Pipeline and independent cross-checks}
The diagonalisation is done in Julia by \fn{HSF/generate\_spectral\_data.jl}, an
optimised successor of the earlier \fn{generate\_EE\_data.jl}, which is left
untouched. The analysis is done in Python, in \fn{code/qca\_fragmentation/spectral/}.
\begin{itemize}
\item \textbf{Two independent constructions of $U$} agree. The Julia operator
  is built from Kronecker products of local gates. The Python operator comes
  from this repository's exact $\mathbb{Z}[1/\sqrt2]$ engine
  (\fn{core/cycle.py}). Their quasienergy multisets agree to $<10^{-14}$ for all
  16 rules at $N=8$ (test \fn{test\_engine\_matches\_julia\_spectrum}, $N=6,7,8$).
  Comparisons of angle multisets cut the circle in the middle of the largest
  gap. Sorting raw angles is unsafe because a level at $\theta=0$ computed as
  $-10^{-16}$ wraps to $2\pi$. A first version of the test fell into exactly
  this trap and reported a spurious disagreement.
\item \textbf{Against the stored HSF eigendata} (complex \fn{eigen} per sector):
  quasienergies agree to $10^{-14}$. The entropies agree to $10^{-15}$ for every
  globally non-degenerate level, the only case where the old entropies are
  basis independent.
\end{itemize}

\subsection{Diagonalisation (details and optimisation in Sec.~\ref{sec:algorithm})}
$U$ is real orthogonal (Sec.~\ref{sec:lit}). Each block is diagonalised by
\begin{enumerate}
\item solving the real symmetric problem $C=(A+A^T)/2$, whose eigenvalues are
  $\cos\theta$, with LAPACK's MRRR solver ($O(n)$ workspace);
\item grouping the $\cos\theta$ into clusters (tolerance $10^{-9}$) and
  diagonalising the antisymmetric part $(A-A^T)/2$ inside each cluster, which
  separates $\theta$ from $-\theta$.
\end{enumerate}
This yields exactly orthonormal eigenvectors, also inside degenerate eigenspaces,
which a general complex eigensolver does not guarantee. It is also several times
faster. The three dense products the cluster loop needs are done as single
matrix--matrix multiplications rather than per cluster (that step was memory-bound
and took half the run time). On the 12th-gen i7 used here, dense LAPACK is
fastest with 4--6 BLAS threads and degrades sharply at 20. Every run stores the
largest eigen-residual $\max|U v-e^{-i\theta}v|$ and the orthonormality error.

\subsection{Symmetry resolution}
\label{sec:symres}
Each eigenvector carries three labels, all found and verified at run time:
\begin{enumerate}
\item \textbf{Krylov sector}: connected component of the support of $U$
  (classical fragmentation).
\item \textbf{Gate-algebra block}: the decomposition of each sector by the
  commutant of the bond algebra $\mathcal{A}_0=\langle g_1,\ldots,g_N\rangle$
  generated by the local gates~\cite{moudgalya2022}. This includes quantum
  fragmentation and hidden internal symmetries. It is found without prior
  knowledge of the symmetry:
  \begin{enumerate}
  \item diagonalise a generic element $K=\sum_ia_ig_i$ (random $a_i\in[1,2]$,
    fixed seed);
  \item transform a second generic element $K'=\sum_ib_ig_i$ into that
    eigenbasis;
  \item take the connected components of the transformed $K'$ as the blocks.
  \end{enumerate}
  If the commutant is Abelian, each eigenvector of $K$ lies in one block, and
  $K'$ connects everything within a block. Isomorphic blocks (non-Abelian
  multiplicity) are merged, which is harmless: they give exactly degenerate
  copies of one spectrum, removed in the statistics. \emph{Check}: the largest
  matrix element of $Q^TUQ$ between different blocks is stored
  (\fn{max\_offblock}).
\item \textbf{Spatial label.}
  \begin{itemize}
  \item obc0, odd $N$, mirror-symmetric rules ($r_{01}=r_{10}$): the reflection $R$ commutes with $U$, and every
    eigenvector gets a parity $\pm1$, a true quantum number.
  \item obc0, even $N$: $R$ exchanges the two layers ($RL_AR=L_B$), so $R$ alone is no symmetry. The half step
    $\Omega=RL_A$ is one: it commutes with $U$ and $\Omega^2=U$ (both asserted at run time).
  \item pbc (even $N$ only), every rule: translation by \emph{one} site maps layer $A$ onto layer $B$, so the
    half-translation $\Omega_T=T_1L_A$ commutes with $U$ and $\Omega_T^2=UT_2$, with $T_2$ the translation by two sites
    (asserted). It resolves the $T_2$ momentum $j=0,\ldots,N/2-1$ and splits each momentum sector once more. The
    site-centred reflection does not commute with $\Omega_T$ in general and is not resolved on the ring: a residual
    $\mathbb{Z}_2$ can remain in the $j=0$ and $j=N/4$ sectors of mirror-symmetric rules.
  \end{itemize}
  For the square-root symmetries ($\Omega$, $\Omega_T$) the stored quantity is their eigenphase $\varphi$, obtained by
  diagonalising the symmetry inside each degenerate eigenspace of $U$ (with the same orthonormal $\cos$/$\sin$
  two-step as for $U$). The resolved spectrum of a block is the spectrum of $\varphi$, not of $\theta$: a $\pm1$
  ``sheet'' label derived from $e^{i\theta/2}$ would flip at the branch cut $\theta=\pm\pi$ and mix the two halves
  there. Labels are used for a block only if the block is checked to be invariant under the symmetry, and the invariance
  residual is stored.
\end{enumerate}
Time reversal ($T=L_AK$, $T^2=1$) fixes the reference ensemble to $\beta=1$
(Sec.~\ref{sec:lit}).

\subsection{Statistics}
Gap ratios are computed per fully resolved block (sector, gate-algebra block, true symmetry label), on the block's level
variable ($\varphi$ where a square-root symmetry was resolved, otherwise $\theta$), with exact degeneracies removed.
\begin{itemize}
\item \emph{Folding is decided numerically per block.} If the block spectrum is symmetric under $x\to-x$ (complex
  conjugation maps the block to itself: every block of a real symmetry), only $|x|\in[0,\pi]$ is used. Otherwise, e.g.\
  in a ring-momentum sector $j\ne-j$ whose conjugate partner is another sector, the full circle is used, including the
  wrap-around spacing.
\item \emph{The reference ensemble depends on the block.} Self-conjugate blocks carry the time reversal $T=L_AK$ and are
  compared with $\beta=1$ ($0.5307$). In momentum sectors $j\ne-j$, $T$ maps $j\to-j$, so no antiunitary symmetry is
  left within the sector, and a chaotic sector is $\beta=2$ (CUE, $0.5996$). Both branches are tested against random
  matrices: CUE through the full-circle branch gives $0.597\pm0.005$, and products of random reflections through the
  folding branch give $0.534\pm0.007$.
\item Only blocks with at least 18 ratios are used, and the ratios of all such blocks are pooled. For comparison the
  same statistic is computed at two coarser resolutions: per Krylov sector only, and for the whole spectrum as one list.
\end{itemize}

\subsection{Entanglement and the fair reference}
For every eigenvector the half-chain von Neumann entropy $S_\mathrm{vN}$ is
stored. For every degenerate eigenspace of dimension $g$ (inside its block and
label) the basis-independent Haar-averaged R\'enyi-2 entropy is stored as well,
\begin{equation}
  S_2^\mathrm{Haar}=-\ln\frac{\operatorname{Tr}P_A^2+\operatorname{Tr}P_B^2}{g(g+1)},
  \qquad P_A=\operatorname{Tr}_BP,\ P_B=\operatorname{Tr}_AP,
  \label{eq:s2haar}
\end{equation}
where $P$ is the projector onto the eigenspace. This follows from
$\mathbb{E}[\ket\psi\!\bra\psi^{\otimes2}]=P^{\otimes2}(1+\mathrm{SWAP})/(g(g+1))$,
and for $g=1$ it is the ordinary R\'enyi-2 entropy of the eigenvector. The same
formula with $P$ the projector onto the whole \emph{block} gives
$S_2^\mathrm{rand}(\mathrm{block})$, the entanglement of a random state inside
that block. That is the fair reference for a fragmented system
(Sec.~\ref{sec:lit}), and the ratio
\begin{equation}
  \eta=S_2^\mathrm{Haar}/S_2^\mathrm{rand}(\mathrm{block})
\end{equation}
is $\approx1$ for a thermal eigenstate and $\ll1$ for an atypically weakly
entangled, scar-like one. $\eta$ is defined for blocks of dimension $\ge2$.
States with $\eta<0.5$ in blocks of dimension $\ge20$ are counted as
scar-like.

% R32: the algorithm and the optimisation choices (input by R32_floquet_spectra.tex)
\section{Algorithm and optimisation choices}
\label{sec:algorithm}

This section documents the production code, \fn{HSF/generate\_spectral\_data\_v2.jl} (Julia), step by step. For every
design choice it gives the reason, the measured effect, and how the change was validated. The original pipeline
(\fn{generate\_EE\_data.jl}, untouched) and the first rewrite (v1, \fn{generate\_spectral\_data.jl}) serve as references.
All timings are on one 12th-generation Intel i7-12700H (6 performance + 8 efficiency cores, 20 hardware threads, 31~GB),
with OpenBLAS.

\subsection{The algorithm}
For one unit (rule $r$, size $N$, gate $V\in\{H,X\}$, boundary condition obc0 or pbc):
\begin{enumerate}
\item \textbf{Local gates from bit operations.} Each controlled gate $g_i$ is a sparse $2^N\times2^N$ matrix with at most
  two entries per column. It is built by reading the two neighbour bits of each basis state (vacuum outside the chain
  for obc0, wrap-around for pbc). For obc0 every $g_i$ is asserted equal to HSF's Kronecker-product construction.
\item \textbf{The Floquet operator is never formed.} $L_A=\prod_{i\ \mathrm{odd}}g_i$ and $L_B=\prod_{i\ \mathrm{even}}g_i$
  are sparse (at most $2^{N/2}$ entries per column). $U$ is applied as $Uv=L_B(L_Av)$.
\item \textbf{Krylov sectors.} For $V=H$: union-find over the supports of the individual gates. This is the same
  partition as the connected components of $\mathrm{supp}(U)$ (R-T19 Prop.~4), asserted at run time for every $N\le10$.
  For $V=X$, $U$ is a permutation and is formed directly; its orbits are the sectors.
\item \textbf{Spatial symmetry}, detected and verified by random-vector tests through the layers:
  \begin{itemize}
  \item obc0: the reflection $R$ if $RUR=U$; otherwise the half step $\Omega=RL_A$, if $RL_AR=L_B$ (asserting $\Omega^2=U$
    and $[\Omega,U]=0$).
  \item pbc: translation $T_2$ by two sites, and the site-centred reflection if it commutes with $U$.
  \end{itemize}
\item \textbf{Per Krylov sector $S$} ($d=|S|$):
  \begin{enumerate}
  \item \emph{Gate-algebra blocks} (Sec.~\ref{sec:symres}): diagonalise $K=\sum_ia_ig_i|_S$ (random $a_i$, fixed seed),
    transform $K'=\sum_ib_ig_i|_S$ into that eigenbasis $Q$, and take connected components. Skipped for $V=X$, whose
    orbits are finer than the gate-algebra blocks, and for $d\le2$.
  \item \emph{Dense sector block} $U_S$, assembled from the layers in column chunks. Then $Q^TU_SQ$, and each block
    $A_b=Q_b^TU_SQ_b$. Check: the largest off-block element is stored.
  \item \emph{Per block}: real-symmetric eigensolve of $C=(A_b+A_b^T)/2$; clustering of $\cos\theta$; within each
    cluster, the Hermitian matrix $iV_c^TS_bV_c$, with $S_b=(A_b-A_b^T)/2$, separates $\theta$ from
    $-\theta$. Then quasienergy
    groups (degeneracies); symmetry labels by successive diagonalisation inside each degenerate eigenspace;
    entanglement.
  \end{enumerate}
\item \textbf{Output} (JLD2, readable with h5py): per eigenvector the sector, block, symmetry label, $\theta$,
  $S_\mathrm{vN}$, degeneracy, $S_2^\mathrm{Haar}$ of its eigenspace, and $S_2^\mathrm{rand}$ of its block. Per unit: the
  sector and block sizes, the symmetry, and five checks (eigen-residual, orthonormality, off-block leakage, symmetry
  residual, run time).
\end{enumerate}

\subsection{Optimisation choices, with measured effect}

\paragraph{1. Real-symmetric eigensolver instead of a complex non-Hermitian one.} $U$ is real orthogonal, so
$C=(U+U^T)/2$ is real symmetric with eigenvalues $\cos\theta$, and it commutes with $S=(U-U^T)/2$.
\begin{itemize}
\item \emph{Speed:} at $n=2048$ (4 BLAS threads) the complex \fn{eigen(U)} takes 11.0~s, a real non-symmetric
  \fn{eigen(U)} 4.4~s, and the real-symmetric \fn{eigen!(Symmetric(C))} 2.2~s, a factor 5.
\item \emph{Correctness:} a general eigensolver does not return orthonormal eigenvectors inside degenerate eigenspaces,
  and these spectra are highly degenerate (up to 95\,\% of the states for 105/150). The symmetric solver does, so every
  stored entanglement belongs to a properly normalised, mutually orthogonal basis.
\item \emph{Splitting:} $\theta$ and $-\theta$ share $\cos\theta$. Inside each $\cos$-cluster, $S$ restricted to the
  cluster is antisymmetric, and $iS$ is Hermitian with eigenvalues $\pm\sin\theta$. Clusters are tiny (mostly pairs),
  so this costs nothing.
\item \emph{Workspace:} the MRRR driver needs $O(n)$ workspace, against $O(n^2)$ for divide and conquer, which matters
  at $n=16384$.
\end{itemize}

\paragraph{2. BLAS threading.} Dense LAPACK on this hybrid CPU is fastest with 4--6 BLAS threads. Measured for a
$4096\times4096$ symmetric eigensolve:
\begin{center}\small
\begin{tabular}{lrrrr}
\toprule
BLAS threads & 1 & 4 & 10 & 20 \\
\midrule
eigensolve (s, Julia $-$t~1) & 9.7 & 5.1 & 6.4 & 32.9 \\
eigensolve (s, Julia $-$t~20) & 11.6 & 5.7 & 6.6 & 7.0 \\
\bottomrule
\end{tabular}
\end{center}
Production uses 6 BLAS threads for the dense stages and Julia's 20 threads (with BLAS set to 1) for the per-cluster loop.

\paragraph{3. Batched products instead of per-cluster ones.} The first version multiplied the full $d\times d$ matrices
by each cluster's few columns separately. With about $d/2$ clusters, this streams the full matrices from memory
$d/2$ times and is memory-bound: 39 of 74~s at $N=12$ (rule 147). Doing $AV$, $SV$ and $EV$ as matrix--matrix products
cut the run to 35~s.

\paragraph{4. Never forming $U$} (the memory incident). Profiling the phases at $N=12$ showed 1.1~GB live after
\emph{building $U$ and finding its sectors}, before any dense algebra. For the chaotic rules $U$ is effectively dense: a
sparse matrix with up to $2^N$ entries per column, plus Graphs.jl's adjacency lists for the component search. That
extrapolates to about 18~GB at $N=14$, which is what triggered the memory reaper. Applying $U$ through the two sparse
layers, and finding sectors from the individual gates, removes this entirely. Peak memory for rule 147 at $N=13$: 7.4~GB
(v1) $\to$ 4.5~GB (v2); run time 224~s $\to$ 124~s.

\paragraph{5. Shared lookup instead of sparse slicing.} \fn{SparseArrays}' \fn{getindex} allocates an $m$-length cache
per call. With $N$ gates sliced per sector and $10^3$--$10^4$ sectors at $N\ge16$, this alone would cost terabytes of
memory traffic. All restrictions now use one shared position vector, set and reset per sector.

\paragraph{6. Entanglement in compressed coordinates.} The entanglement of an eigenvector supported on a sector $S$ only
involves the rows and columns of the $2^{N_B}\times2^{N_A}$ reshaped matrix that $S$ occupies. Zero rows and columns
change neither the singular values nor $\operatorname{Tr}P_A^2$ and $\operatorname{Tr}P_B^2$, so the code works with the
occupied rows and columns only. For fragmented rules this shrinks every SVD from $256\times256$ (at $N=16$) to a few
dozen. Effect for rule 156: $N=16$ from $>10$~minutes (aborted) to 15~s; $N=18$ in 36~s with 1.3~GB peak. These are what
make the per-rule larger sizes possible.

\paragraph{7. Smaller choices.}
\begin{itemize}
\item The orthonormality check uses 64 sampled columns for blocks larger than 2048 (the full $V^TV$ would be another
  $d\times d$ matrix).
\item Garbage collection is forced only after blocks with $d>2000$. Forcing it after every block made rule 156 at $N=9$
  ten times slower (2.2~s $\to$ 21~s), since it has hundreds of tiny blocks.
\item Block bases are stored as contiguous copies, not views. A view with non-contiguous columns made $E\,V$ fall back
  from BLAS to Julia's generic single-threaded product (observed: about 1.2 cores busy).
\item $Q^TXQ$ is formed in column chunks. This did \emph{not} lower the measured peak (4.5~GB either way): the peak sits
  in the block eigensolve itself (block, basis, eigenvectors, LAPACK workspace).
\item Julia runs with \fn{--heap-size-hint=14G}, so the collector works harder before the machine runs short.
\end{itemize}
Expected peak at $N=14$ (a $16384$-dimensional block): about 15~GB.

\subsection{How every change was validated}
No optimisation was accepted on the strength of its own checks. Each one was compared against the previous version, and
the whole chain against independent constructions:
\begin{center}\small
\begin{tabular}{p{0.36\textwidth}p{0.2\textwidth}p{0.36\textwidth}}
\toprule
comparison & scope & result \\
\midrule
v1 vs HSF eigendata (complex \fn{eigen} per sector) & rules 4, 6, 15, $N=9$ & $\theta$ to $10^{-14}$; $S_\mathrm{vN}$ (non-degenerate) to $10^{-15}$ \\
exact Python engine vs Julia & all 16 rules, $N=6,7,8$ & spectra to $<10^{-9}$ (test suite) \\
v2 vs v1, after each change & 6--8 rules, $N=9$ & block by block identical$^\dagger$ \\
v2 (pbc) vs exact engine (pbc) & 7 rules, $N=8$ & spectra to $<10^{-14}$ \\
run-time assertions & every run & gates $=$ kron (obc0); sectors $=$ components of $\mathrm{supp}(U)$ ($N\le10$); $\Omega^2=U$, $[\Omega,U]=0$, $[T_2,U]=0$ \\
\bottomrule
\end{tabular}
\end{center}
$^\dagger$ Compared as multisets per block: sector size, block size, sorted quasienergies, $S_2^\mathrm{Haar}$,
degeneracies, $S_\mathrm{vN}$ of the non-degenerate states, and $S_2^\mathrm{rand}$.

\subsection{Pitfalls caught, recorded so they are not repeated}
\begin{enumerate}
\item \textbf{Shared variables in threaded closures (Julia).} Twice, a variable assigned inside a
  \fn{Threads.@threads} loop had the same name as one in the enclosing function (\fn{SV}, then \fn{G}). Julia then
  shares one boxed variable across threads, a race that produced dimension errors. Fixed by renaming. Any new variable in
  \fn{process\_block!} must not reuse a name assigned in the loop body.
\item \textbf{Comparing angles on a circle.} Sorting raw angles misaligns lists when a level at $\theta=0$ is computed as
  $-10^{-16}$ and wraps to $2\pi$. A first comparison therefore reported a spurious 0.7~rad disagreement and briefly
  stopped production. All comparisons now cut the circle in the largest gap (\fn{floquet.circle\_distance}).
\item \textbf{Sort-key comparisons.} Pairing states by a sort on (sector, block, $\theta$ rounded) misaligns nearly
  degenerate levels and reported a false mismatch between two correct versions. All comparisons are now block-wise
  multisets.
\item \textbf{Shell tooling.} \fn{pkill -f} with a pattern that also appears in one's own command line kills that
  command. A \fn{for} loop over paths split them at the spaces in ``QCA Simulations'' (caught after one minute; the
  deletions targeted non-existent fragments, so no file was lost).
\item \textbf{An incomplete symmetry analysis on the ring.} The first pbc runs resolved only $T_2$ and folded every
  block to $|\theta|$. The chaotic rules then came out at $\langle r\rangle\approx0.43$ on the ring against $0.53$ with
  open boundaries, a typical two-sector mixture. Two errors were responsible. First, the half-translation
  $\Omega_T=T_1L_A$ commutes with $U$ for every rule on an even ring and was not resolved. Second, a momentum sector
  $j\ne-j$ is not conjugation-symmetric, so folding it pairs unrelated levels. Both were fixed (Sec.~\ref{sec:symres}),
  the 64 affected files regenerated (the old ones kept in \fn{results\_spectra/H\_pbc\_T2only/}), and the
  statistics given a numerical fold test. The same analysis showed that the $\pm1$ label of $\Omega$ (obc0, even $N$)
  flips at the branch cut. Its blocks were therefore regenerated with the eigenphase $\varphi$ stored.
\item \textbf{The memory reaper.} Background jobs that push the machine to critically low memory are stopped by the
  session and must not be restarted without permission. This was the motivation for items~4--7.
\end{enumerate}

% R32 section 4: results, gate H, obc0 (input by R32_floquet_spectra.tex)
\section{Results: the sixteen rules with the Hadamard gate}
\label{sec:results}

\subsection{The organising principle: how many neighbour patterns fire the gate}

Count, for each rule, the number $k$ of the four neighbour patterns $(x_{i-1},x_{i+1})\in\{00,01,10,11\}$ on which the
Hadamard acts. For $k=2$, also record whether the firing set depends on only one neighbour (``one-sided'': $60,102,153,195$)
or on both (``XOR'': $150$ fires on $01,10$, and $105$ on $00,11$). This single combinatorial number sorts the 16 rules into
classes with distinct spectral signatures (Table~\ref{tab:summary}, Fig.~\ref{fig:overview}):

\begin{center}
\small
\begin{tabular}{llp{0.5\textwidth}}
\toprule
class & rules & spectral signature (obc0) \\
\midrule
$k=0$ & $204$ & $U=I$ (identity) \\
$k=1$ & $108,156,198,201$ & exponential fragmentation; Poisson inside the pooled blocks \\
$k=2$ one-sided & $60,102,153,195$ & a tower of gate-algebra blocks; $\langle r\rangle$ near Poisson, drifting up with $N$.
  \emph{On a ring: no tower, chaotic} (Sec.~\ref{sec:pbc}) \\
$k=2$ XOR & $105,150$ & arithmetic spectrum, $3^m$ distinct quasienergies (W150) \\
$k=3$ & $54,57,99,147$ & \textbf{chaotic (GOE)}, plus a fixed number of exact gate-singlet eigenstates.
  On a ring: chaotic, CUE or COE by symmetry \\
$k=4$ & $51$ & $U=H^{\otimes N}$, non-interacting \\
\bottomrule
\end{tabular}
\end{center}

\begin{table}[t]
\centering
\scriptsize
\IfFileExists{tab_r32_summary_H.tex}{\resizebox{\textwidth}{!}{\begin{tabular}{llrlrrrrrrrr}
\toprule
$k$ & $W$ & $N$ & sym. & sectors & blocks & $\langle r\rangle_\mathrm{none}$ & $\langle r\rangle_\mathrm{full}$ & \#$r$ & distinct $\theta/2^N$ & $\bar\eta_\mathrm{big}$ & singlets \\
\midrule
0 & 204 & 16 & R & 65536 & 65536 & -- & -- & 0 & 1.53e-05 & -- & 0 \\
1 & 108 & 19 & R & 55405 & 124345 & 0.386 & 0.392 & 71048 & 0.0303 & 0.55 & 52688 \\
1 & 156 & 19 & none & 10946 & 121370 & 0.412 & 0.421 & 5486 & 0.00408 & 0.36 & 37812 \\
1 & 198 & 19 & none & 10946 & 121370 & 0.412 & 0.421 & 5486 & 0.00408 & 0.35 & 37812 \\
1 & 201 & 19 & R & 31572 & 90053 & 0.383 & 0.389 & 97266 & 0.068 & 0.54 & 39872 \\
2 one-sided & 60 & 14 & none & 15 & 105 & 0.384 & 0.395 & 7732 & 0.492 & 0.68 & 24 \\
2 one-sided & 102 & 14 & none & 15 & 105 & 0.384 & 0.395 & 7732 & 0.492 & 0.68 & 24 \\
2 one-sided & 153 & 14 & none & 1 & 15 & 0.398 & 0.421 & 8092 & 0.992 & 0.81 & 2 \\
2 one-sided & 195 & 14 & none & 1 & 15 & 0.398 & 0.421 & 8092 & 0.992 & 0.81 & 2 \\
2 XOR & 105 & 15 & R & 9 & 17 & 0.491 & 0.498 & 1512 & 0.0186 & 0.81 & 7 \\
2 XOR & 150 & 15 & R & 9 & 23 & 0.404 & 0.389 & 5987 & 0.1 & 0.66 & 7 \\
3 & 54 & 14 & Omega & 2 & 3 & 0.428 & 0.523 & 8127 & 0.992 & 0.84 & 1 \\
3 & 57 & 14 & none & 1 & 3 & 0.531 & 0.531 & 8190 & 1 & 0.96 & 2 \\
3 & 99 & 14 & none & 1 & 3 & 0.531 & 0.531 & 8190 & 1 & 0.96 & 2 \\
3 & 147 & 14 & Omega & 1 & 5 & 0.422 & 0.527 & 8117 & 0.991 & 0.89 & 4 \\
4 & 51 & 14 & R & 1 & 16384 & -- & -- & 0 & 0.000122 & -- & 16384 \\
\bottomrule
\end{tabular}
}}{\fbox{table pending}}
\caption{Per-rule summary at each rule's largest computed $N$ (obc0, gate $H$), generated by
\fn{spectral.analysis.write\_table}. sym.: spatial symmetry resolved ($R$ for odd $N$; the half-step $\Omega$ for even $N$).
sectors/blocks: Krylov sectors, and fully resolved gate-algebra blocks. $\langle r\rangle_\mathrm{none}$: whole spectrum as one
list. $\langle r\rangle_\mathrm{full}$: pooled over fully resolved blocks with $\ge20$ distinct half-levels, from \#$r$ ratios.
distinct $\theta/2^N$: number of distinct quasienergies over the whole spectrum. $\bar\eta_\mathrm{big}$: mean
$S_2/S_2^\mathrm{rand}(\mathrm{block})$ in the largest block. singlets: one-dimensional gate-algebra blocks inside Krylov
sectors of dimension $\ge3$, i.e.\ exact eigenstates of every individual gate.}
\label{tab:summary}
\end{table}

\begin{figure}[t]
\centering
\ifgraphic{../../figures/r32/fig_r32_overview_H.pdf}{\textwidth}
\caption{(a) Gap ratio of every rule at its largest $N$: unresolved (whole spectrum, open circles) and fully resolved
(Sec.~\ref{sec:symres}, filled). Lines: Poisson, GOE ($\beta=1$, the class of $U$), GUE. (b) Fraction of distinct
quasienergies in the whole spectrum (log scale). Rules grouped by the number $k$ of firing neighbour patterns.}
\label{fig:overview}
\end{figure}

\subsection{$k=3$: chaos, with exact scars that do not grow}
With open boundaries, the four rules on which the gate fires on all but one pattern are the only chaotic ones. (On a
ring, the one-sided $k=2$ rules join them; Sec.~\ref{sec:pbc}.)
\begin{itemize}
\item \textbf{57 and 99} (mirror images; computed independently, identical to all digits): $\langle r\rangle=0.5311\pm0.0028$
  at $N=14$ (8190 ratios), on the GOE value $0.5307$ from $N=10$ on: $0.525, 0.523, 0.527, 0.535, 0.531$ for
  $N=10\ldots14$. $P(r)$ follows the $\beta=1$
  surmise over the whole range (Fig.~\ref{fig:rule57}). The eigenstates in the big block are thermal \emph{within that
  block}: $\bar\eta=0.95$.
\item \textbf{147}: $\langle r\rangle_\mathrm{full}=0.526$ ($N=13$) and $0.5275\pm0.0028$ ($N=14$), but only after resolving
  the reflection (odd $N$) or the half step $\Omega$ (even $N$). Unresolved, the same spectra give $0.424$ and $0.422$. This is the textbook symmetry-mixing effect,
  and here it is the difference between ``integrable'' and ``chaotic''.
\item \textbf{54}: $\langle r\rangle_\mathrm{full}$ rises steadily, $0.441, 0.456, 0.502, 0.510, 0.523$ for
  $N=10\ldots14$, towards GOE. Its frozen vacuum $|0\ldots0\rangle$ is a separate Krylov sector.
\end{itemize}

\paragraph{Gate singlets.} Inside the chaotic block, each of these rules carries one-dimensional blocks of the gate
algebra: states that are eigenvectors of \emph{every individual gate} $g_i$, hence of $U$ for any order and any number of
gate applications. Their number does not depend on $N$ ($N=6\ldots13$):
\begin{itemize}
\item two for 57 and 99, and four for 147, all \emph{product} states ($S=0$) at $\theta\in\{0,\pi\}$;
\item exactly one for 54, at $\theta=0$, which is \emph{entangled}. Its entanglement \emph{decreases} with $N$:
  $S_\mathrm{vN}=0.470, 0.430, 0.405, 0.371, 0.345, 0.315, 0.291, 0.265, 0.243$ for $N=6\ldots14$ ($S_2$ from 0.35 to
  0.13). It is an area-law state that approaches a product state as the chain grows. (An earlier draft of this report said
  ``0.29 at every $N$''; that generalised from the single value at $N=12$ and is withdrawn.)
\end{itemize}
These are exact many-body scars in the sense of Refs.~\cite{shiraishi2017,moudgalya2022review}: non-thermal eigenstates
embedded in a GOE spectrum. The embedding is exact here, by the gate algebra, not approximate. Because they are
eigenstates of every gate, they are also stationary under any reordering of the circuit, which a Hamiltonian scar cannot
offer.

\paragraph{N-stable outliers in 57/99.} Besides the singlets, 57 has a handful of anomalously weakly entangled eigenstates
in the big block ($\eta\approx0.2$): 4 at $N=11$ and $13$, and 10 at $N=12$. Their quasienergies are stable with $N$
($\theta/\pi=\pm0.6877\to\pm0.6879$ from $N=11$ to $13$, and $\pm0.3122$ at $N=12$ and $13$). They do not form an equally
spaced tower, and their number does not grow. This points to a localised structure (a single defect in a frozen
background) rather than a PXP-type scar band. In 147, by contrast, the low-$\eta$ states are numerous (72--83) and spread
over $\theta$: the finite-size tail of the thermal distribution.

\subsection{$k=1$: exponential fragmentation; two pairs at the same $N$}
Rules 108, 156, 198 and 201 fire on a single pattern (108 on $11$, 201 on $00$, 156 on $10$, 198 on $01$) and have
exponentially many Krylov sectors (R2, R19). They were run to $N=20$ (201 to $N=19$; see below). The class is compared
at the largest $N$ that all four rules share, $N=19$ (Table~\ref{tab:summary}), and not each rule at its own maximum.
\begin{itemize}
\item \textbf{108 and 201} (fire on $11$ or $00$; mirror symmetric) have large blocks: the largest has 2583 and 5544
  states at $N=19$. Pooled over them, $\langle r\rangle=0.3922\pm0.0010$ (108) and $0.3891\pm0.0009$ (201), with
  $7\times10^4$ to $10^5$ ratios. This is Poisson ($0.3863$) up to a small, statistically resolved excess of
  0.003 to 0.006, which is stable from $N=17$ on.
\item \textbf{156 and 198} (fire on $10$ or $01$; a reflection pair, computed independently and identical) stay
  \emph{tiny}: the largest block has 100 states at $N=19$ (125 at $N=20$), and only $2139$ of the $2^{19}$ quasienergies are
  distinct. $\langle r\rangle=0.4209\pm0.0038$ at $N=19$ and $0.4172\pm0.0023$ at $N=20$, stable from $N=17$ on.
  This is clearly \emph{above} Poisson but far below GOE.
\end{itemize}
The block structure explains both values. The commutant study (R-T18/R-T27) shows that a block is a tensor product of
independent rooms, so its quasienergies are sums of room quasienergies. When many rooms of different sizes contribute
(108, 201), the sums are effectively independent and the statistics are Poisson. When a block is built from a few small,
repeated rooms (156, 198), the sums are \emph{structured}. Exact coincidences are removed, but (our reading, not tested) the surviving levels keep
an arithmetic regularity that suppresses small spacings. Neither pair shows level repulsion growing with $N$.

The gate singlets are exponentially many (at $N=19$: 52\,688 for 108, 37\,812 for 156 and 198, and 39\,872 for 201;
roughly doubling with each added site). These are the dark states of the rooms.

\subsection{$k=2$, XOR: arithmetic spectra}
150 and 105 have only a handful of Krylov sectors (the domain-wall shells, R8) and large blocks, but their spectra are
\emph{arithmetic}. The number of distinct quasienergies of 150 is exactly $3^m$ for even $N$ and $(3^m+(-1)^m)/2$ for odd
$N$ ($m=\lceil N/2\rceil$) at every $N=6\ldots14$, reproducing R8's law from an independent diagonalisation. 105 shares
the even-$N$ counts but has far fewer at odd $N$ (251 against 1093 at $N=13$). Up to 95\,\% of the eigenstates lie in
degenerate eigenspaces, and spacings repeat exactly. $\langle r\rangle$ (0.42 and 0.36) is therefore no measure of chaos
here. The right description is the free wall walk with sign defects of R8. Each Krylov sector contains one gate singlet at
$\theta=0$, an entangled dark state ($S$ between 0.13 and 1.14).

\subsection{$k=2$, one-sided: a boundary-driven tower, statistics undecided}
In 60/102 (gate fires iff the left or right neighbour is 1) and 153/195 (iff it is 0), the gate on a site is controlled by
one neighbour only. With open boundaries the gate algebra then produces a tower:
\begin{itemize}
\item 153/195: $N+1$ blocks of sizes $2^{N-1},2^{N-2},\ldots,2,1$, plus two singlets;
\item 60/102: $2N-4$ singlets and many small blocks.
\end{itemize}
At $N=5$ the dominant basis states of block $j$ (153) are those whose first site from the left in the $-1$ eigenstate of
$H$ is site $j$. The sites to its left are frozen in the dressed $+1$ state, and the right part is a smaller copy of the
chain. The entanglement accordingly falls into bands, one per block (Fig.~\ref{fig:rule153}). Within the blocks,
$\langle r\rangle$ is near Poisson but \emph{increases} steadily with $N$: for 60, $0.334, 0.351, 0.360, 0.369, 0.381,
0.395$ for $N=9\ldots14$; for 153, $0.37$--$0.414$ for $N=9\ldots13$. Integrability and a slow crossover towards GOE cannot
be told apart at these sizes \emph{on the open chain}. The tower itself is a boundary effect. On a ring it is absent, and
the same rules are chaotic with $\beta=1$ statistics ($\langle r\rangle=0.531$ for 153 at $N=14$; Sec.~\ref{sec:pbc}). So
on the open chain, the non-ergodicity of this class is produced by the ends.

\subsection{$k=0$ and $k=4$: trivial}
204 is the identity (one distinct quasienergy). In 51 every site gets $H$ in its layer, so $U=H^{\otimes N}$ exactly. Its
spectrum is $\{0,\pi\}$, every state is a gate singlet, and the eigenstates are products of single-site $H$ eigenstates,
with $S=0$.

\begin{figure}[p]
\centering
\ifgraphic{../../figures/r32/fig_r32_rule_H_W57.pdf}{\textwidth}
\caption{Rule 57 ($k=3$), largest $N$, obc0. (a)~Half-chain entanglement against quasienergy. Coloured: non-degenerate states,
by block. Grey: states in a degenerate eigenspace, whose entanglement depends on the basis. Diamonds: gate singlets (exact
eigenstates of every gate; product states at $\theta=0,\pi$). (b)~$\eta=S_2^\mathrm{Haar}/S_2^\mathrm{rand}(\mathrm{block})$,
basis independent. (c)~$P(r)$ of the fully resolved half spectrum against the Poisson, GOE and GUE surmises. (d)~$\langle
r\rangle$ against $N$ at three resolution levels.}
\label{fig:rule57}
\end{figure}

\begin{figure}[p]
\centering
\ifgraphic{../../figures/r32/fig_r32_rule_H_W147.pdf}{\textwidth}
\caption{Rule 147 ($k=3$), largest $N$; panels as in Fig.~\ref{fig:rule57}. In (d) the Krylov-only statistics (here the whole
spectrum) stay near Poisson, while resolving the spatial symmetry ($R$ for odd $N$, the half step $\Omega$ for even $N$)
brings $\langle r\rangle$ to the GOE value.}
\label{fig:rule147}
\end{figure}

\begin{figure}[p]
\centering
\ifgraphic{../../figures/r32/fig_r32_rule_H_W153.pdf}{\textwidth}
\caption{Rule 153 ($k=2$, one-sided), largest $N$, obc0; panels as in Fig.~\ref{fig:rule57}. The entanglement bands in (a) are the
blocks of the boundary-driven tower (sizes $2^{N-1},2^{N-2},\ldots$). $P(r)$ is close to Poisson, while $\langle r\rangle$
drifts upward with $N$.}
\label{fig:rule153}
\end{figure}


\section{Periodic boundary conditions}
\label{sec:pbc}
% R32 section 5: periodic boundary conditions (input by R32_floquet_spectra.tex)
% Numbers come from the generated tables tab_r32_bc_H_N14.tex and tab_r32_shape_H_pbc_N14.tex
% (qca_fragmentation.spectral.analysis.write_bc_table / write_shape_table).
The user asked whether the classes survive a change of boundary condition. The obvious suspects are the one-sided
rules, whose block tower is anchored at the open ends. All sixteen rules were run on rings of even length $N=6\ldots14$. At $N=14$, the eight rules with $k\ge2$ that are not
arithmetic have one $\approx16\,000$-dimensional sector each, and each run took about 16~min.

\subsection{What is resolved on the ring, and which ensemble to expect}
On an even ring, translation by one site $T_1$ maps layer $A$ onto layer $B$, for every rule. The half-step
\begin{equation}
  \Omega_T = T_1 L_A,\qquad [\Omega_T,U]=0,\qquad \Omega_T^2 = U\,T_2,
\end{equation}
is therefore a symmetry of every rule. Its eigenphase $\varphi$ fixes both the quasienergy and the momentum $j$ of the
two-site translation $T_2$ ($e^{2i\varphi}=e^{-i\theta}e^{2\pi i j/(N/2)}$). A statistics block is
(Krylov sector, gate-algebra block, $j$), and its level variable is $\varphi$. The site reflection $R$ does \emph{not}
commute with $\Omega_T$ and is not resolved. Which random-matrix ensemble applies depends on the type of sector:
\begin{itemize}
\item \emph{Real sectors} ($j=0$, and $j=N/4$ when $N/4$ is an integer) are mapped to themselves by complex conjugation
  $K$. At $j=0$, $T_1$ is an involution, so $\Omega_T$ is again a product of two reflections: $\beta=1$ on the half
  spectrum (rule~57: $0.549$ at $N=14$). At $j=N/4$, $T_1^2=-1$ and the reflection structure is lost, so the reference is
  a real orthogonal matrix without that structure, which has $\beta=2$-like local statistics (rule~57: $0.610$ at
  $N=12$). For mirror-symmetric rules, the real sectors also contain the unresolved $\mathbb Z_2$ generated by $L_AR$,
  which pushes $\langle r\rangle$ down. \textbf{The real sectors are therefore not used as the diagnostic.}
\item \emph{Complex sectors} ($j\ne -j$) carry most of the levels, so they are the diagnostic. $K$ maps $j\to-j$, so a
  time-reversal symmetry inside sector $j$ needs a unitary $P$ that also flips $j$ and satisfies $PUP^{-1}=U^{-1}$. For a
  mirror-symmetric rule the site reflection does this ($R$ exchanges the layers), so the naive reference is COE ($\beta=1$,
  $0.5307$). For a rule without mirror symmetry the naive reference is CUE ($\beta=2$, $0.5996$).
\end{itemize}
The block classification (real or complex) is numerical: a block counts as real if its eigenphases are closed under
conjugation to $10^{-7}$.

\paragraph{Correction during the analysis.} The first version of this conjugation test compared the sorted lists of $x$
and $-x$ element by element. A level at $0\equiv2\pi$ (the $j=0$ sector has forced eigenvalues $\pm1$) moves to the
other end of the sorted list and shifts every pair. The test then wrongly classified $j=0$ as complex, so its full circle
was used. The fix (\fn{stats.conj\_symmetric}) matches each $-x$ to its nearest level on the circle. A test pins the
$j=0$ block of rule~57 at $N=8$ against the exact Python engine. The obc0 numbers of Sec.~\ref{sec:results} were
recomputed: none of the 145 units changed.

\subsection{Results}
\begin{table}[h]
\centering\small
\resizebox{\textwidth}{!}{\begin{tabular}{llrrrrrrrlrr}
\toprule
 & & \multicolumn{2}{c}{sectors} & \multicolumn{2}{c}{blocks} & obc & \multicolumn{3}{c}{pbc $\langle r\rangle$} & \multicolumn{2}{c}{singlets} \\
$k$ & $W$ & obc & pbc & obc & pbc & $\langle r\rangle_\mathrm{full}$ & real $j$ & complex $j$ & naive ref. & obc & pbc \\
\midrule
0 & 204 & 16384 & 16384 & 16384 & 16384 & -- & -- & -- & COE & 0 & 0 \\
1 & 108 & 3329 & 2627 & 5694 & 4517 & 0.383 & 0.404 & 0.386 & COE & 2070 & 1667 \\
1 & 156 & 987 & 844 & 5740 & 5072 & -- & -- & -- & CUE & 2439 & 2016 \\
1 & 198 & 987 & 844 & 5740 & 5072 & -- & -- & -- & CUE & 2439 & 2016 \\
1 & 201 & 1897 & 2627 & 4143 & 4517 & 0.395 & 0.403 & 0.452 & COE & 1760 & 1667 \\
2 one-sided & 60 & 15 & 2 & 105 & 9 & 0.395 & 0.414 & 0.511 & CUE & 24 & 1 \\
2 one-sided & 102 & 15 & 2 & 105 & 9 & 0.395 & 0.414 & 0.511 & CUE & 24 & 1 \\
2 one-sided & 153 & 1 & 2 & 15 & 9 & 0.421 & 0.428 & 0.531 & CUE & 2 & 1 \\
2 one-sided & 195 & 1 & 2 & 15 & 9 & 0.421 & 0.427 & 0.531 & CUE & 2 & 1 \\
2 XOR & 105 & 8 & 7 & 15 & 20 & 0.356 & 0.405 & 0.391 & COE & 7 & 7 \\
2 XOR & 150 & 8 & 10 & 15 & 52 & 0.419 & 0.424 & 0.440 & COE & 7 & 6 \\
3 & 54 & 2 & 2 & 3 & 9 & 0.523 & 0.383 & 0.483 & COE & 1 & 1 \\
3 & 57 & 1 & 1 & 3 & 8 & 0.531 & 0.549 & 0.602 & CUE & 2 & 1 \\
3 & 99 & 1 & 1 & 3 & 8 & 0.531 & 0.549 & 0.602 & CUE & 2 & 1 \\
3 & 147 & 1 & 2 & 5 & 9 & 0.527 & 0.424 & 0.518 & COE & 4 & 1 \\
4 & 51 & 1 & 1 & 16384 & 16384 & -- & -- & -- & COE & 16384 & 16384 \\
\bottomrule
\end{tabular}
}
\caption{Open chain against ring at $N=14$, all sixteen rules. ``naive ref.''\ is
the expectation for the complex sectors from the identified symmetries (mirror $\Rightarrow$ COE). The one-sided rules
violate it (Table~\ref{tab:shape}). The $N=12$ table is \fn{tab\_r32\_bc\_H\_N12.tex}.}
\label{tab:bc}
\end{table}

\begin{table}[h]
\centering\small
\begin{tabular}{llrrrrrrrr}
\toprule
$k$ & $W$ & mirror & \#$r$ & $\langle r\rangle$ & $P(r<0.15)$ & $\chi^2_\mathrm{P}$ & $\chi^2_{\beta=1}$ & $\chi^2_{\beta=2}$ & $\chi^2_\mathrm{P+\beta=2}$ \\
\midrule
3 & 57 & no & 13944 & 0.602 & 0.023 & 8482 & 1083 & 47 & 47 \\
3 & 99 & no & 13944 & 0.602 & 0.023 & 8482 & 1083 & 47 & 47 \\
3 & 147 & yes & 13938 & 0.518 & 0.082 & 3444 & 124 & 4324 & 400 \\
3 & 54 & yes & 13938 & 0.483 & 0.105 & 2119 & 850 & 9869 & 466 \\
2 one-sided & 60 & no & 13938 & 0.511 & 0.079 & 3228 & 163 & 4873 & 399 \\
2 one-sided & 102 & no & 13938 & 0.511 & 0.079 & 3228 & 163 & 4873 & 399 \\
2 one-sided & 153 & no & 13938 & 0.531 & 0.063 & 4191 & 27 & 2458 & 385 \\
2 one-sided & 195 & no & 13938 & 0.531 & 0.063 & 4191 & 27 & 2458 & 385 \\
\bottomrule
\end{tabular}

\caption{Shape of $P(r)$ in the complex momentum sectors of the ring, $N=14$: $\chi^2$ of a 10-bin histogram (9 degrees
of freedom) against the Atas surmises, and against the best Poisson+CUE mixture. A crossover from Poisson towards CUE
can pass through the COE \emph{mean}, but then its small-$r$ tail is too heavy ($P(r)\sim r^\beta$). The surmise values
of $P(r<0.15)$ are 0.261 (Poisson), 0.064 ($\beta=1$) and 0.019 ($\beta=2$). The surmise deviates from large-$N$ random
matrices by about 1\%, so $\chi^2$ of a few tens at $n\approx1.4\times10^4$ is a good fit.}
\label{tab:shape}
\end{table}

\begin{figure}[h]
\centering
\ifgraphic{../../figures/r32/fig_r32_bc_compare_H.pdf}{\textwidth}
\caption{Open chain against ring, each rule at the largest even $N$ available for both boundaries ($N=14$ for the
$k\ge2$ non-arithmetic rules, $N=12$ otherwise). (a) Resolved $\langle r\rangle$. For the ring, the real momentum sectors
(open squares) and the complex sectors (filled squares) are shown separately. (b) Number of resolved blocks.}
\label{fig:bc}
\end{figure}

\paragraph{$k=3$: the class survives, and the ensemble follows the symmetry.}
57 and 99 have no mirror symmetry and become CUE on the ring: $\langle r\rangle=0.602\pm0.002$, with a GUE-shaped
$P(r)$ ($\chi^2=47$, against 1083 for $\beta=1$). This is the textbook outcome for a chaotic system once its momentum
sectors are complex. 147 (mirror) is COE-shaped and approaches $0.531$ from below ($0.469$, $0.496$, $0.518$ for
$N=10,12,14$). 54 (mirror) follows more slowly ($0.458$, $0.472$, $0.483$). On the open chain, 54 also lagged the other
three rules by about two sizes.

\paragraph{One-sided rules: the tower is a boundary effect, and on the ring they are chaotic.}
On the ring the one-sided rules have 2 Krylov sectors and 9 blocks, against a tower of $N+1$ blocks (60: 105 blocks at
$N=14$). In the complex sectors their statistics are $\beta=1$ in \emph{shape}:
\begin{itemize}
\item 153 and 195: $\langle r\rangle=0.531$, $P(r<0.15)=0.063$, and $\chi^2_{\beta=1}=27$ against 385 for the best
  Poisson+CUE mixture.
\item 60 and 102: still rising ($0.423$, $0.471$, $0.511$), and already closer to COE than to any mixture.
\end{itemize}
This contradicts the naive reference: these rules are \emph{not} mirror symmetric, so one expects CUE. There must
therefore be an antiunitary symmetry inside each momentum sector that we have not identified. A search at $N=8$ found nothing beyond the known layer swap $T_1UT_1^{-1}=U^{-1}$, which preserves $j$ and so gives no
time reversal inside a sector. It looked for $W$ with $WUW^{-1}=U^{\pm1}$, and separately with
$W\Omega_TW^{-1}=\Omega_T^{-1}$, over $W=D\,L\,R^aT_s$: $a\in\{0,1\}$, all $s$, $L$ the identity or a uniform or
sublattice product of $X$ or $H$ (and of $Z$ without $D$), and $D$ an \emph{arbitrary} diagonal phase operator, solved
exactly as a consistency problem on the support graph. For 147 the search returns the reflection, which validates it.
The symmetry is an open point
(Sec.~\ref{sec:open-pbc}).

The consequence for the classification is that \textbf{the one-sided class is defined by its boundary}. With open ends
the tower splits the space into $N+1$ blocks and the statistics stay near Poisson. On the ring the same bulk dynamics is
chaotic.

\paragraph{$k=2$ XOR: still arithmetic.} The ring makes the spectra of 150 and 105 \emph{more} degenerate. At $N=14$,
150 has 255 distinct quasienergies and 105 has 251, against $2187=3^7$ for both on the open chain. Their statistics
are not those of a chaotic spectrum either: $\langle r\rangle=0.440$ (150) and $0.391$ (105) in the complex sectors,
with a small-$r$ tail of $P(r<0.15)=0.22$ and $0.45$ (Poisson: 0.26; $\beta=1$: 0.06). That tail is the signature of
near-coincident levels in an arithmetic spectrum (Table~\ref{tab:shape}). The spectra are arithmetic, not chaotic, as on
the open chain.

\paragraph{$k=1$: still exponentially fragmented.} At $N=14$ the ring has 2627 sectors for 108 and 201 (open chain:
3329 and 1897) and 844 for 156 and 198 (open: 987). The blocks stay small: the largest has 143 states for 108/201 and 25
for 156/198. On the ring, 108 and 201 have \emph{identical} block structures (sectors, blocks, singlets and the count
of distinct levels), but different statistics ($0.399$ and $0.417$, from 2547 ratios each). This is consistent with a
relabelling of the basis that maps one rule's sector graph onto the other's, while the gate $H$ is not invariant under
it. The same holds for 156 and 198, which are mirror images. Within the large uncertainty of such small blocks, the
statistics stay between Poisson and the structured value of the open chain.

\paragraph{Exact singlets on the ring.} Every $k=3$ and one-sided rule has exactly \emph{one} gate singlet on the ring, at
$\theta=0$, for every $N$ (the XOR rules keep about one per sector, as on the open chain). The singlets at $\theta=\pi$ of the open chain are gone, so they are boundary states. For
57, 99, 147, 153 and 195 the ring singlet is a product state ($S_\mathrm{vN}=0$). For 54, and also for 60 and 102, it is
entangled, with the \emph{same} entropy as the entangled singlet of 54 on the open chain:
$S_\mathrm{vN}=0.4045,\ 0.3447,\ 0.2908,\ 0.2430$ for $N=8,10,12,14$. Equal entropies across two rules and two boundary
conditions suggest one and the same state. The eigenvectors were not stored, so overlaps could not be checked.

\subsection{What survives}
\label{sec:open-pbc}
\begin{center}
\begin{tabular}{lp{0.3\textwidth}p{0.42\textwidth}}
\toprule
class & open chain & ring \\
\midrule
$k=3$ & chaotic (COE) & chaotic: CUE (57, 99) or COE (147, 54), as the symmetry dictates \\
$k=2$ one-sided & boundary tower, near Poisson & tower gone, chaotic ($\beta=1$ shape) \\
$k=2$ XOR & arithmetic & arithmetic, even more degenerate \\
$k=1$ & exponential fragmentation & unchanged \\
$k=0,4$ & trivial & trivial \\
\bottomrule
\end{tabular}
\end{center}
The $k=3$, XOR, $k=1$ and trivial classes are robust. The one-sided class is not a bulk class: its non-ergodicity comes
from the open ends. ``Chaos needs three firing patterns'' is therefore a statement about the open chain. On the ring,
two firing patterns suffice if they depend on a single neighbour.

\paragraph{Open points.}
\begin{enumerate}
\item The antiunitary symmetry behind the $\beta=1$ statistics of the one-sided rules on the ring.
\item Resolving $L_AR$ in the real sectors of the mirror rules, which needs the stored eigenvectors or one extra
  symmetry in the Julia pipeline.
\item The identity of the entangled singlet shared by 54 (both boundaries) and 60/102 (ring).
\end{enumerate}


\section{The X gate}
\label{sec:x}
% R32 section 6: the X gate, compared with H (input by R32_floquet_spectra.tex)
% Numbers: tab_r32_X_N14.tex (spectral.xgate.write_x_table), results_spectra/summary_X*.json,
% results_spectra/X_cycle_growth.json (spectral.xgate.cycle_growth).
With $V=X$ every rule is a classical reversible cellular automaton. $U$ permutes the basis, so its spectrum is fixed
by the cycle structure alone: a cycle of length $L$ contributes the $L$-th roots of unity, once each. The eigenvectors
are the Fourier modes of the cycles,
\[
  |q,c\rangle = L^{-1/2}\sum_{t=0}^{L-1} e^{2\pi i q t/L}\,U^t|x_c\rangle .
\]
Level statistics therefore measure orbit arithmetic, not chaos, and are not reported. The informative quantities are
the cycle-length distribution and the entanglement of the Fourier modes. A mode is supported on $L$ basis states, so
its Schmidt rank is at most $L$ and $S_\mathrm{vN}\le\ln L$: short cycles force low entanglement.

\paragraph{Data and checks.} All sixteen rules were run with the same Julia pipeline, on the open chain for $N=6\ldots14$
and on rings for even $N=6\ldots14$: 224 runs and 40~minutes in total, with no failure. For every run, the stored
quasienergies equal the union of the roots of unity over the cycles to $2\times10^{-14}$. The cycles were computed
independently in Python from the permutation (\fn{spectral.xgate.cycle\_lengths}). A test checks this against the dense
permutation matrix at $N=8$, for both boundaries. The cycle growth in Fig.~\ref{fig:xcycles} was computed in Python up
to $N=18$.

\subsection{What the sixteen automata are, classically}
Writing $m,n$ for the left and right neighbour, the X-gate update is $x_i\to x_i\oplus f(m,n)$, with $f$ the firing
pattern:
\begin{itemize}
\item \textbf{Linear or affine} ($f$ of degree one): the XOR rules 150 ($f=m\oplus n$) and 105 ($1\oplus m\oplus n$),
  and all one-sided rules (60: $f=m$; 102: $f=n$; 153: $1\oplus n$; 195: $1\oplus m$). These are solvable by linear
  algebra over $\mathbb F_2$. X-102 (flip iff the right neighbour is 1) is exactly the deterministic Floquet quantum East
  model, whose quench dynamics are solved exactly~\cite{bertini2024det}; 60 is its mirror image.
\item \textbf{Known interacting integrable automata.} 54 fires iff at least one neighbour is 1: this is RCA54, the
  interacting integrable Rule~54 with solitons~\cite{buca2021}. 201 fires iff both neighbours are 0: this is the
  Floquet-PXP automaton~\cite{wilkinson2020}. On the ring the global spin flip is a symmetry of the X automata. It maps
  147 onto 54, 108 onto 201, 156 onto 198 and 57 onto 99. With $H$ it is not a symmetry, and the obc0 boundary (ghost sites
  fixed to 0) breaks it too: for example, 108 and 201 have 4868 and 3492 cycles on the open chain at $N=14$, but 4156
  each on the ring.
\item \textbf{Trivial}: 204 ($U=1$) and 51 ($U=X^{\otimes N}$).
\item \textbf{Nonlinear and not known to be integrable}: only 57 and 99 (fire unless $(m,n)=(0,1)$, or $(1,0)$).
\end{itemize}

\subsection{Cycle structure: only 57/99 have exponentially long orbits}
\begin{figure}[h]
\centering
\ifgraphic{../../figures/r32/fig_r32_X_cycles.pdf}{\textwidth}
\caption{X gate: (a) the longest cycle and (b) the mean cycle length seen by a uniformly random basis state,
$\langle L\rangle_\mathrm{state}=\sum_c L_c^2/2^N$, over $2^N$. Filled: open chain; open symbols: ring. Computed from the
permutation up to $N=18$.}
\label{fig:xcycles}
\end{figure}
On the open chain, the longest cycle of 57/99 holds a \emph{finite fraction} of the Hilbert space: $129\,880$ of $2^{18}$
states at $N=18$. A random state lies on a cycle of average length $0.22$ to $0.50$ times $2^N$ for every $N$ from 6 to 18.
Classically, this is an ergodic orbit structure. Every other rule has orbits that grow at most linearly:
\begin{itemize}
\item RCA54: longest cycle $3N-5$ for 54 (147, which is RCA54 only on the ring: 26 at $N=14$);
\item XOR: $N+1$;
\item one-sided: 16 or 32, independent of $N$ (a period that is a power of two, as for linear automata);
\item Floquet-PXP (201, 108): slow growth, 187 states at $N=18$.
\end{itemize}
On the ring, the cycles of 57/99 still grow exponentially, but more slowly (4428 states at $N=18$), and their share of the
space falls to $\langle L\rangle_\mathrm{state}/2^N=0.007$ at $N=14$. The classically ergodic orbit of 57/99 on the open
chain needs the boundary. The reason is a conservation law: on the ring X-57 conserves the number of domain walls, which the
fixed-$0$ ends of the open chain create and destroy (Sec.~\ref{sec:answers}). This is the mirror image of the one-sided rules with $H$, which need the ring to become chaotic.

\subsection{Eigenstate entanglement: H against X at one N}
\begin{table}[h]
\centering\scriptsize
\begin{tabular}{llrrrrr|rrrrr}
\toprule
 & & \multicolumn{5}{c|}{open chain} & \multicolumn{5}{c}{ring} \\
$k$ & $W$ & cycles & longest & $\langle L\rangle/2^N$ & $\bar S_X$ & $\bar S_H$ & cycles & longest & $\langle L\rangle/2^N$ & $\bar S_X$ & $\bar S_H$ \\
\midrule
0 & 204 & 16384 & 1 & 6.1e-05 & 0.00 & 0.00 & 16384 & 1 & 6.1e-05 & 0.00 & 0.00 \\
1 & 108 & 4868 & 33 & 4.3e-04 & 0.11 & 0.06 & 4156 & 77 & 5.3e-04 & 0.17 & 0.12 \\
1 & 156 & 2280 & 60 & 7.9e-04 & 0.14 & 0.02 & 2202 & 35 & 6.5e-04 & 0.19 & 0.04 \\
1 & 198 & 2280 & 60 & 7.9e-04 & 0.14 & 0.02 & 2202 & 35 & 6.5e-04 & 0.19 & 0.04 \\
1 & 201 & 3492 & 55 & 6.9e-04 & 0.14 & 0.07 & 4156 & 77 & 5.3e-04 & 0.17 & 0.11 \\
2 one-sided & 60 & 1044 & 16 & 9.7e-04 & 0.44 & 0.47 & 300 & 63 & 3.8e-03 & 0.84 & 0.89 \\
2 one-sided & 102 & 1044 & 16 & 9.7e-04 & 0.44 & 0.47 & 300 & 63 & 3.8e-03 & 0.84 & 0.89 \\
2 one-sided & 153 & 1024 & 16 & 9.8e-04 & 0.47 & 0.66 & 300 & 63 & 3.8e-03 & 0.84 & 0.92 \\
2 one-sided & 195 & 1024 & 16 & 9.8e-04 & 0.47 & 0.66 & 300 & 63 & 3.8e-03 & 0.84 & 0.92 \\
2 XOR & 105 & 1096 & 15 & 9.1e-04 & 0.59 & 0.52 & 1172 & 14 & 8.5e-04 & 0.58 & 0.75 \\
2 XOR & 150 & 1096 & 15 & 9.1e-04 & 0.59 & 0.51 & 2344 & 7 & 4.3e-04 & 0.43 & 0.57 \\
3 & 54 & 632 & 37 & 1.6e-03 & 0.71 & 0.89 & 732 & 91 & 2.1e-03 & 0.65 & 0.82 \\
3 & 57 & 34 & 4288 & 2.4e-01 & 0.98 & 0.98 & 500 & 462 & 7.0e-03 & 0.69 & 0.99 \\
3 & 99 & 34 & 4288 & 2.4e-01 & 0.98 & 0.98 & 500 & 462 & 7.0e-03 & 0.69 & 0.99 \\
3 & 147 & 632 & 26 & 1.6e-03 & 0.70 & 0.93 & 732 & 91 & 2.1e-03 & 0.65 & 0.93 \\
4 & 51 & 8192 & 2 & 1.2e-04 & 0.16 & 0.00 & 8192 & 2 & 1.2e-04 & 0.16 & 0.00 \\
\bottomrule
\end{tabular}

\caption{X gate at $N=14$, both boundaries: number of cycles, the longest, $\langle L\rangle_\mathrm{state}/2^N$, and the
mean half-chain entanglement of all eigenstates over the Page value ($\bar S_X$). $\bar S_H$ is the same mean for the H
gate, over its non-degenerate eigenstates (degenerate ones have basis-dependent entanglement).}
\label{tab:x}
\end{table}

\begin{figure}[h]
\centering
\ifgraphic{../../figures/r32/fig_r32_HX_compare_N14.pdf}{\textwidth}
\caption{H against X at $N=14$, open chain. (a) Mean eigenstate entanglement over the Page value. (b) X gate:
$\langle L\rangle_\mathrm{state}/2^N$.}
\label{fig:hx}
\end{figure}

Table~\ref{tab:x} and Fig.~\ref{fig:hx} give the following picture:
\begin{itemize}
\item \textbf{57/99:} on the open chain, X and H are \emph{equally} thermal ($0.98$ of Page). The cycle Fourier modes of an
  ergodic orbit look like random states. On the ring, the X value drops to $0.69$ with the shorter cycles, while H stays
  at $0.99$.
\item \textbf{54/147 (RCA54):} X gives $0.70$, bounded by its short orbits ($\ln 37=3.6$ against Page 4.35). H gives
  $0.89$ and $0.93$, and its spectrum is GOE (Sec.~\ref{sec:results}). \textbf{The Hadamard turns an interacting integrable
  classical automaton into a chaotic quantum circuit.}
\item \textbf{One-sided:} $0.44$ to $0.47$ (X) against $0.47$ to $0.66$ (H) on the open chain. On the ring both rise
  (X: $0.84$; H: $0.89$ to $0.92$), but only H is chaotic there; X stays affine, with period at most a few hundred.
\item \textbf{XOR:} X ($0.59$) is slightly \emph{above} H ($0.52$) on the open chain; both are arithmetic.
\item \textbf{$k=1$:} both are low (X $0.11$ to $0.14$; H $0.02$ to $0.07$). These are fragmented rooms, whatever the gate.
\end{itemize}

\subsection{What the comparison says}
\begin{enumerate}
\item Quantum chaos with $H$ does \emph{not} require classical chaos with $X$. Two of the four chaotic $k=3$ rules, 54 and
  147, are the interacting integrable RCA54 with $X$, yet GOE-chaotic with $H$. On the ring, 57/99 keeps its CUE chaos
  with $H$, although its X orbits no longer fill a finite share of the space. What the chaotic rules share is the firing
  count $k=3$, not the classical dynamics. Deforming classical automata by non-commuting gates is the setting of
  Ref.~\cite{iadecola2020}; for these rules, the deformation by $H$ yields chaos.
\item Conversely, the only classically ergodic automaton (57/99 on the open chain) is also chaotic with $H$.
\item With $X$, every rule splits into cycles, so ``fragmentation'' is total and carries no information. With $H$, the
  fragmentation is a property of the rule class: exponential at $k=1$ on both boundaries, and boundary-driven for the
  one-sided rules.
\end{enumerate}


% R32 section: thermalisation beyond the gap ratio (2026-10-15).  Input by R32_floquet_spectra.tex.
% Numbers: tab_r32_thermal.tex (spectral.thermal.write_thermal_table), results_spectra/thermal_*.{json,npz}.
\section{Thermalisation beyond the gap ratio}
\label{sec:thermal}
The gap ratio is the most robust single diagnostic, but it probes only nearest-neighbour level correlations of whole
symmetry blocks. Three complementary diagnostics were added:
\begin{itemize}
\item the \emph{spectral form factor} (SFF), which measures long-range spectral correlations and the time scale of
  ergodicity;
\item \emph{ETH for a local observable}, i.e.\ whether the individual eigenstates are thermal;
\item \emph{quench dynamics} from product states, i.e.\ whether specific initial states thermalise.
\end{itemize}
Thermal means infinite temperature inside the block a state lives in.

\begin{table}[h]
\centering\small
\begin{tabular}{llrrrrrrr}
\toprule
 & & ring $\langle r\rangle$ & ETH & $t_2$ & \multicolumn{2}{c}{random product} & \multicolumn{2}{c}{scar $+\,|h_-\rangle$} \\
$k$ & $W$ & ($N$) & exponent & ($N=16$) & open & ring & open & ring \\
\midrule
3 & 57 & 0.598 (16) & -0.40 & 78 & 1.00 & 1.00 & 0.95 & 0.96 \\
3 & 54 & 0.510 (16) & -0.13 & 990 & 0.99 & 0.99 & 0.31 & 0.44 \\
3 & 147 & 0.528 (16) & -0.27 & 426 & 1.00 & 1.00 & 0.63 & 0.60 \\
2 one-sided & 153 & 0.527 (16) & -0.24 & 564 & 0.91 & 1.00 & 0.00 & 0.89 \\
2 one-sided & 60 & 0.523 (16) & -0.17 & 747 & 0.72 & 0.99 & 0.00 & 0.78 \\
2 XOR & 150 & 0.431 (14) & -- & -- & 0.56 & 0.69 & 0.26 & 0.42 \\
1 & 201 & 0.452 (14) & -- & -- & 0.08 & 0.14 & 0.29 & 0.31 \\
1 & 156 & -- & -- & -- & 0.00 & 0.01 & 0.11 & 0.14 \\
\bottomrule
\end{tabular}

\caption[Thermalisation summary.]{Thermalisation summary. Ring $\langle r\rangle$: complex momentum sectors at the
largest $N$ (in brackets). ETH exponent: slope of $\log\sigma$ against $\log d$ for the eigenstate-to-eigenstate
fluctuation $\sigma$ of the magnetisation density $m=\sum_iZ_i/N$ in the complex ring sectors ($N=8\ldots14$); ETH
predicts $-1/2$. It is shown only for rules whose momentum sectors are single Krylov sectors. $t_2$: time after which
the log-binned SFF stays within a factor 2 of the RMT ramp, at $N=16$. Last four columns: $S_\mathrm{half}/S_\mathrm{Page}$
averaged over $t=200\ldots300$ after a quench at $N=14$, from random computational product states (average of 8) and
from $|h_+\rangle^{\otimes N}$ with one $|h_-\rangle$ in the middle.}
\label{tab:thermal}
\end{table}

\subsection{Spectral form factor}
The SFF $K(t)=\langle|\mathrm{Tr}_bU^t|^2\rangle/d$ was computed over the complex momentum sectors on the ring. There each
block's eigenphases fill the circle and the reference is COE or CUE (Sec.~\ref{sec:pbc}). On the open chain every block is
real, with a $\pm\theta$-symmetric spectrum, so the standard form factors do not apply. A single circuit has no ensemble,
and the SFF does not self-average, so the curves were averaged over the sectors and over logarithmic time bins. At
$t\gtrsim100$ each bin pools hundreds of points, so the residual noise is a few percent.

Fig.~\ref{fig:ethsff}(b) shows two systematic features:
\begin{enumerate}
\item A non-universal excess at short times ($t\approx5\ldots30$ periods): at $N=16$ about $10^2$ times the ramp for 57
  and 153, and $10^3$ to $10^4$ for 54, 60 and 147. It grows with $N$ for every rule, most strongly for 147 and 60.
\item A slow approach to the ramp. For 57, $K/K_\mathrm{RMT}\approx1.4$--$1.6$ persists over two decades and reaches 1.1
  only near the Heisenberg time $D$.
\end{enumerate}
A sharp Thouless time (agreement within $\pm30\%$) is therefore not reached well before $D$ at these sizes. The robust
proxy $t_2$ ranks the rules clearly (Fig.~\ref{fig:ethsff}(c)). For 57 it grows slowly, from 25 to 78 periods for
$N=10\ldots16$ (roughly as $D^{0.3}$). For the other chaotic rules it is 5 to 13 times longer at $N=16$ and grows faster.
The early excess suggests that part of the space of these rules has approximately periodic dynamics. This is a tentative
reading; it fits the slow scar-defect dynamics below.

\subsection{ETH for a local observable}
The diagonal matrix elements $\langle n|m|n\rangle$ of the magnetisation density were computed in all non-degenerate
eigenstates of the complex ring sectors ($N=8\ldots14$). On the open chain the central $Z$ was used in the largest
gate-algebra block, at even $N$; this gives the same ordering. For 57 the fluctuation $\sigma$ falls as
$d^{-0.40}$ (open chain: $d^{-0.36}$), close to the ETH value $d^{-1/2}$ at these sizes. The other chaotic rules converge
markedly more slowly: $-0.27$ (147), $-0.24$ (153), $-0.17$ (60) and $-0.13$ (54). This is the same ordering as their
$\langle r\rangle$ and SFF convergence. For 150 nearly all eigenstates are degenerate (free fermions), so a diagonal ETH
test is meaningless there. The $k=1$ momentum sectors contain many Krylov sectors, so the test does not apply to them
either.

\begin{figure}[h]
\centering
\ifgraphic{../../figures/r32/fig_r32_eth_sff.pdf}{\textwidth}
\caption{(a) ETH: eigenstate-to-eigenstate fluctuation of the magnetisation density in the complex ring sectors, against
the sector dimension; dashed: $d^{-1/2}$. (b) SFF over the RMT ramp at $N=16$, log-binned. (c) $t_2$ against $N$.
Colours as in (a).}
\label{fig:ethsff}
\end{figure}

\subsection{Quench dynamics}
Product states were evolved with the sparse layers at $N=14$ for 300 periods (Fig.~\ref{fig:dyn},
Table~\ref{tab:thermal}). The findings, class by class:
\begin{itemize}
\item \emph{$k=3$:} random product states reach the Page value on both boundaries (0.99 to 1.00). For 57 so do Néel and
  the scar with one defect.
\item \emph{One-sided, a boundary effect in real time:} random states saturate at 0.72 (60) and 0.91 (153) of Page on
  the open chain, but at 0.99 to 1.00 on the ring.
\item \emph{XOR (150):} stays well below Page (0.56 to 0.69). Néel does not move at all on the ring.
\item \emph{$k=1$:} hardly entangles (201: 0.08 to 0.14, 156: about 0); this is the fragmentation.
\item \emph{The scar:} $|h_+\rangle^{\otimes N}$ returns with probability 1 for every rule.
\item \emph{One $|h_-\rangle$ defect on top of the scar:} slow for 54 and 147. Their entanglement is still rising at
  $t=300$ (0.3 to 0.6 of Page), while for 57 it saturates within about 100 periods. This is a bulk analogue of the edge
  mode of 57: a localised excitation on the inert $|h_+\rangle$ vacuum that leaks into the bulk only slowly. On the open
  chain the one-sided rules keep the defect state at exactly zero half-chain entanglement. Everything to the left of the
  defect is the inert $|h_+\rangle$ part of the tower, so the cut sees a product state.
\end{itemize}

\begin{figure}[h]
\centering
\ifgraphic{../../figures/r32/fig_r32_dynamics_N14.pdf}{\textwidth}
\caption{Half-chain entanglement over the Page value after quenches, $N=14$. (a) Random computational product states,
open chain (average of 8); (b) the same on the ring; (c) $|h_+\rangle^{\otimes N}$ with one $|h_-\rangle$ in the middle,
for the $k=3$ rules (solid: open chain, dotted: ring).}
\label{fig:dyn}
\end{figure}

\paragraph{Verdict.} The ordering of all diagnostics is the same: 57/99 are the fastest and cleanest thermalisers
(CUE/COE statistics, ETH-like fluctuations, short $t_2$). Behind them come 153, 147, 60 and 54, all chaotic on the
ring, with the convergence slowing in that order. XOR and $k=1$ do not thermalise. The non-thermal exceptions inside the
chaotic rules are all tied to the inert $|h_+\rangle$ vacuum: the exact scars, the edge mode of 57, and the slow defect
states of 54 and 147.


% R32 section: follow-up (2026-10-15) -- the open points, answered.  Input by R32_floquet_spectra.tex.
% Every statement here is pinned by a test in code/qca_fragmentation/tests/test_spectral.py.
\section{Exact structures behind the classes}
\label{sec:answers}
This section collects the exact structures found in the follow-up (2026-10-15), which explain the classes: the flip map,
the hidden time reversal of the one-sided rules, the universal scar, the kink dispersion of 156/198 and the domain-wall
charge of X-57. It also reports the block-first runs beyond the dense-sector limit (step~9), the eigenvector analyses at
small $N$ (step~10), and the checks of the assumptions those runs make. Every statement is pinned by a test in
\fn{code/qca\_fragmentation/tests/test\_spectral.py}.

\subsection{Flip partners: the X flip turns $H$ into $XHX$}
The physical spin flip is $X^{\otimes N}$. It maps a rule $W$ onto its partner $W'$ (the firing pattern
$f(m,n)\to f(1-m,1-n)$), but the gate changes with it:
\begin{equation}
  X^{\otimes N}\,U_W(H)\,X^{\otimes N}=U_{W'}(XHX),\qquad XHX=\tfrac{1}{\sqrt2}(X-Z)=-ZHZ .
\end{equation}
Because $Z^{\otimes N}$ commutes with the controls, $Z^{\otimes N}U_{W'}(G)Z^{\otimes N}=U_{W'}(ZGZ)$. So
$U_W(H)\cong U_{W'}(-H)$ up to the basis change $Z^{\otimes N}$: the partner rule with a $-1$ on every firing. This is a
conditional phase, not a global one. Verified as operators at $N=8$, on the ring, for the pairs 54/147, 108/201, 60/195,
102/153 and 57/99. The consequences are:
\begin{itemize}
\item On the ring, flip partners have \emph{identical} supports, hence identical Krylov sectors and blocks, but
  \emph{different} spectra. Examples: 108 and 201 ($\langle r\rangle=0.399$ and $0.417$ from identical block structures),
  60/102 and 153/195 (0.511 and 0.531), and 54 and 147 (0.483 and 0.518).
\item The obc0 boundary (ghost sites fixed to $0$) breaks the flip already for $X$. With $X$ on the ring, 108 and 201 have
  4156 cycles each; on the open chain they have 4868 and 3492.
\item 57/99 and 156/198 are also \emph{reflection} partners, and the reflection is an exact symmetry with either gate.
  Their spectra are identical in every dataset.
\end{itemize}

\subsection{The hidden symmetry of the one-sided rules}
On the ring, the one-sided rules have $\beta=1$ statistics without mirror symmetry (Sec.~\ref{sec:pbc}). The
time-reversal symmetry behind this is now identified:
\begin{equation}
  S=B_e\,D\,R_\mathrm{bond}\,B_e^{\mathsf T},\qquad A=K\,L_A\,S,
\end{equation}
where:
\begin{itemize}
\item $R_\mathrm{bond}$ is the bond-centred reflection $i\to1-i$, which exchanges the two sublattices;
\item $B_e$ rotates the even sites into the eigenbasis $\{|h_+\rangle,|h_-\rangle\}$ of the Hadamard;
\item $D=X^{\otimes N}$ for 153/195, and a diagonal sign operator for 60/102.
\end{itemize}
Verified as explicit operators at $N=8$ and $10$, for all four rules, to $10^{-13}$:
\begin{itemize}
\item $S$ is real, $[S,U]=0$, and $ST_2S^{-1}=T_2^{-1}$ ($S$ flips the momentum);
\item $A$ preserves the momentum and maps $U\to U^{-1}$ and $\Omega_T\to\Omega_T^{-1}$;
\item $A^2=+1$: COE, not Kramers pairs.
\end{itemize}
In the mixed frame (Hadamard basis on the targets of one layer, computational basis on its controls) that layer is
diagonal, and $S$ becomes a dressed reflection. This is why a search in the computational basis could not find it. The
single-site rotation needed is $V=(|h_+\rangle,|h_-\rangle)$, a rotation by $\pi/8$, not $H$. For 57/99 the same search
returns nothing, consistent with CUE.

\paragraph{Correction to Sec.~\ref{sec:pbc}.} In the real momentum sectors ($j=0$, $N/4$), $S$ preserves $j$ and commutes
with $U$. So the real sectors of the one-sided rules carry an unresolved $\mathbb Z_2$, exactly as the reflection does for
the mirror rules. Their real-sector values (0.414 and 0.428 at $N=14$) mix two symmetry sectors and are not
diagnostic. The conclusions used only the complex sectors, so they are unchanged.

\subsection{Scars: one universal eigenstate, edge-projected}
\textbf{$|h_+\rangle^{\otimes N}$ is a $+1$ eigenvector of every controlled-Hadamard gate, for every rule and boundary},
because $H|h_+\rangle=|h_+\rangle$ whatever the controls read. All gate singlets of the chaotic rules follow from it
(found from the gate algebra at $N\le12$ and identified site by site):
\begin{itemize}
\item \emph{$\theta=0$:} $|h_+\rangle^{\otimes N}$ itself.
\item \emph{Rules with a frozen vacuum} (54 on both boundaries, 60/102 on the ring): its projection off the vacuum,
  $|\psi\rangle\propto|h_+\rangle^{\otimes N}-\cos^N(\pi/8)\,|0\ldots0\rangle$. This is the ``entangled singlet''. It has
  Schmidt rank 2, its entanglement follows in closed form and matches the data to $10^{-6}$ for $N=6\ldots14$, and it
  falls to zero as $\cos^N(\pi/8)\to0$. It is the same state in all three rules (overlap $1$).
\item \emph{Boundary singlets} ($\theta=\pi$, or $0$ with two of them): $|h_-\rangle$ on an end site whose gate fires
  for both values of the inner neighbour. With the ghost fixed to $0$, that site picks up $-1$ per period
  deterministically. This predicts the counts exactly: 57 (right end) 2, 99 (left end) 2, 147 (both ends) 4, 54 (neither
  end, vacuum frozen) 1. On the ring there are no ends, hence exactly one singlet.
\end{itemize}
These are exact scars, but of the simplest kind: a product eigenstate of the gate (or its projection) embedded in a
chaotic spectrum, in the spirit of the projector embedding~\cite{shiraishi2017}.

\paragraph{The outliers of 57.} These are \emph{not} singlets. They are an \emph{edge mode on the
$|h_+\rangle^{\otimes N}$ vacuum}: away from the left end the state is $|h_+\rangle$ on every site (overlap rising from
0.86 to 0.93 into the bulk), and the disturbance is localised on the first four sites. The two quasienergies are one
mode: $\pm0.3122\pi$ alone, and $\pm0.6877\pi=\pm0.3122\pi+\pi$ together with the right-end $|h_-\rangle$ singlet. The mode
lives at the left end because there 57's gate depends on the inner neighbour (it fires on $00$, not on $01$); the right
end is inert. On the ring the strongly non-thermal states disappear (minimum $\eta$ 0.41 to 0.55, against 0.13 to 0.27
on the open chain).

\subsection{156/198: one kink per segment, with an exact dispersion}
For 156 (fires only on left $=1$, right $=0$), a $01$ domain wall is frozen and a $10$ wall is mobile. Between two frozen
walls a segment of $L$ free sites holds exactly one kink, in $L+1$ positions. No gate couples two segments, so
$U=\bigotimes_s U_s$. The single-kink spectrum is
\begin{equation}
  \cos\theta_m=-\sin^2\!\frac{\pi m}{L+1},\quad m=1\ldots\lfloor(L+1)/2\rfloor\ (\pm\text{ pairs; }\theta=\pi\text{ at }m=\tfrac{L+1}{2}),
  \qquad\text{plus }\theta=0,
\end{equation}
i.e.\ $\cos\theta(q)=(\cos q-1)/2$ with $q=2\pi m/(L+1)$. It does not depend on the brickwork offset, and it reproduces
the exact spectra of the largest Krylov sectors at $N=12$ as sums over their segments (including the end segments).
Blocks have at most about 125 states, so their levels are sums of a handful of values from this small set. That gives
many exact and near coincidences, hence $\langle r\rangle\approx0.42$: neither independent levels nor level repulsion.

\subsection{X-57: the ergodic orbits need the boundary}
On the ring, X-57 (and X-99) \emph{conserves the number of domain walls}. A search over sublattice-resolved 3- and
4-site densities finds exactly this one local charge (plus its 4-site extension) on the ring, and none on the open
chain. There the ghost sites create and annihilate walls. Ring orbits are confined to fixed wall number, so their share
of $2^N$ decays (0.007 at $N=14$). On the open chain the boundary mixes the wall sectors and the orbits fill a finite
fraction. For contrast, 54 and 150 keep local charges on both boundaries.

\IfFileExists{R32_answers_step9.tex}{% R32: beyond the dense-sector limit (step 9), eigenvectors at small N (step 10), and the checks of their assumptions.
% Data: results_spectra/blockfirst/{ring,tower,verify}_*; code: spectral/blockfirst.py; tests: test_blockfirst_matches_pipeline.
\subsection{Beyond the dense-sector limit (block-first)}
The production pipeline diagonalises whole Krylov sectors densely, so its memory is set by the largest \emph{sector}
(Sec.~\ref{sec:memory}). For the cases that matter for the open questions, an exactly known symmetry was used to split
the space \emph{before} the dense step, by a sparse construction (\fn{spectral/blockfirst.py}), and only eigenvalues were
computed.
\begin{itemize}
\item \emph{Ring, $N=16$} (57, 54, 147, 153, 60): $T_2$-momentum orbit bases first ($\approx8200$ states per sector),
  then the half step $\Omega_T$ inside each sector. The eigenphases come from a Cayley transform, a Hermitian
  eigenvalue problem.
\item \emph{Open chain, one-sided rules:} the largest tower block, built directly in its product basis. For 153 it is
  $|h_-\rangle$ on site 0 with the other sites free ($N=15$); for 60 it is site $0=|1\rangle$ and site $N-1=|h_-\rangle$
  ($N=16$). Each has $16\,384$ states.
\end{itemize}
Validation against the production data: the ring code reproduces the Julia $\varphi$ of every momentum sector to
$10^{-14}$ ($N=12$) and the per-block $\langle r\rangle$ to $10^{-11}$ ($N=12$, 14). The tower blocks reproduce the
Julia largest blocks level by level to $10^{-13}$ ($N=10\ldots13$).

\emph{One pitfall, caught by that comparison:} computing $|\theta|$ from $\arccos$ of the eigenvalues of
$(U+U^{\mathsf T})/2$ smears exactly degenerate eigenvalues $\pm1$ by $\sqrt{\epsilon}\approx10^{-8}$ into spurious
near-duplicates. The real tower blocks therefore use the eigenvalues of $U$ directly.

\emph{Results} (Secs.~\ref{sec:pbc} and \ref{sec:results}):
\begin{itemize}
\item On the ring at $N=16$: 57 is $0.5985\pm0.0011$ (CUE), 147 is $0.528$ and 153 is $0.527$ (COE-shaped), and 60
  ($0.523$) and 54 ($0.510$) are still rising, now COE-shaped.
\item On the open chain the one-sided statistics stop drifting: $\langle r\rangle\approx0.424$ for 153 at $N=15$ and for
  60 at $N=16$.
\end{itemize}

\paragraph{Checking the assumptions (user request).} Block-first assumes that nothing splits the pieces further. This
was tested at the production sizes with the generic-$K$ trick:
\begin{itemize}
\item \emph{Ring:} restricted to a momentum sector, $K$ is a generic combination of $T_2$-invariant elements of the gate
  algebra ($S_A$, $S_B$, $\{S_A,S_B\}$, $U+U^\dagger$). One component per Krylov sector means that no hidden symmetry
  splits the sector: any such symmetry would also split this subalgebra.
\item \emph{Tower:} $K=\sum_ia_ig_i$ restricted to the block.
\end{itemize}
Results:
\begin{itemize}
\item In every tested complex momentum sector at $N=16$ (57 $j{=}1$; 54, 153 $j{=}1,2$; 147 $j{=}1$; 60 $j{=}1,3$) the
  test gives exactly one component, equal to the single Krylov component. $K$ is non-degenerate (gaps
  $4\times10^{-6}$ to $2\times10^{-4}$), and the sector is invariant under $\Omega_T$ to $2\times10^{-15}$.
\item The tower block of 153 at $N=15$ is irreducible: one component of $16\,384$, gap $2.3\times10^{-6}$, and the gates
  leave the block invariant to $1.4\times10^{-14}$.
\item The tower block of 60 at $N=16$ is irreducible: one component of $16\,384$, gap $2.3\times10^{-7}$, and gate
  invariance $2\times10^{-14}$.
\end{itemize}
In the \emph{real} momentum sectors the same test splits the space into two halves. This is how the residual
$\mathbb Z_2$ of the one-sided rules was found, which led to their time reversal (above). The conclusions use only the
complex sectors.

\paragraph{Eigenvectors at small $N$ (step 10).} The scar analyses use eigenvectors from the exact engine at
$N\le12$:
\begin{itemize}
\item the gate singlets, from the one-dimensional gate-algebra blocks;
\item the outliers of 57 and the ETH matrix elements, from the dense eigenvectors of $U$ in the largest block.
\end{itemize}
The Julia pipeline was not changed, so no full recompute was needed.
}{}


\section{Compute and memory}
\label{sec:memory}
% R32: compute log, memory, and the interruptions (numbers from results_spectra/runlog_v2.txt)
\paragraph{Budget and what was used.} The task came with a budget of 8 hours of compute. The user later lifted it
(``finish within the next 12 hours''). All runs used one 12th-generation i7 with 31~GB of memory, one job at a time.
\begin{itemize}
\item First night (v1 pipeline, 01:53 to 02:51): obc0, all 16 rules for $N=6\ldots13$ and $N=14$ for 9 rules; about
  1.7~h including validation and memory tests.
\item Continuation (v2 pipeline, from 12:32 on 2026-10-14): 485 successful runs and 7.1~h of wall time.
  \begin{center}\small
  \begin{tabular}{p{0.6\textwidth}rr}
  \toprule
  stage & runs & wall time \\
  \midrule
  $H$, obc0 (rest of $N=14$; $N=15\ldots20$ by rule; even-$N$ reruns of the mirror rules with $\varphi$) & 117 & 3.55~h \\
  $H$, rings $N=6\ldots14$ (with $\Omega_T$) & 144 & 2.92~h \\
  $X$, obc0 $N=6\ldots14$ & 144 & 0.44~h \\
  $X$, rings $N=6\ldots14$ & 80 & 0.23~h \\
  \bottomrule
  \end{tabular}
  \end{center}
  Every completed run exited with code 0. Runs cut off by the interruptions below were rerun, except rule~153 at
  $N=15$. The 27 non-zero exit codes in the log are the 12:32 start of the runner, which split paths at the spaces in the repository path; those runs
  failed instantly and were rerun after the fix.
\end{itemize}
The most expensive runs are the single $\approx16\,000$-dimensional sectors at $N=14$: 15 to 18~minutes each, with a peak
resident memory of 14~GB.

\paragraph{Interruptions.}
\begin{enumerate}
\item \emph{First night, 02:51}: the session's memory reaper stopped all jobs. The cause was the dense $U$ of the chaotic
  rules: for these rules $U$ is effectively dense, and storing it as a sparse matrix and running a component search
  over it cost 1.1~GB at $N=12$, or about 18~GB at $N=14$. Fixed in v2 (below); resumed with the user's go-ahead.
\item \emph{About 20:08}: the machine shut down. The run lost was 54 at $N=14$, whose old file had already been removed for
  the rerun with $\varphi$. The user then stopped the runs for a second reboot. The rerun was done after the restart.
\item \emph{About 22:45}: the memory reaper stopped all jobs during rule~153 at $N=15$ (above 17~GB). That run is the
  only one that never finished. It was not restarted (see the next paragraph).
\end{enumerate}

\paragraph{Why the sizes stop where they do.} The pipeline diagonalises each Krylov sector densely (the gate-algebra
blocks are resolved inside it), so the memory is set by the largest \emph{sector}. A $16\,383$-dimensional sector peaks at
about 14~GB, and the memory grows faster than the square of the dimension (201: 2.6~GB at $N=18$, sector 6765; 7.7~GB at
$N=19$, sector $10\,946$). The next sizes would need:
\begin{itemize}
\item 153/195 at $N=15$ and 60/102 at $N=16$: a sector of $2^{15}$, about 50~GB;
\item 150/105 at $N=16$: about twice the $12\,870$-dimensional sector of $N=15$, more than 30~GB;
\item 201 at $N=20$: a sector of $17\,711$ ($F_{22}$), about 20 to 23~GB, which is too close to the limit at which the
  session reaper stops every job.
\end{itemize}
These sizes were therefore not attempted, so every class is compared at the largest $N$ that all its members share. That
is $N=14$ for $k=3$ and for the one-sided rules, $N=15$ for XOR, and $N=19$ for $k=1$ on the open chain; on the ring,
$N=14$ for every rule. Going further would need a block-first diagonalisation, which projects onto the gate-algebra
blocks before the dense step. That is the obvious next optimisation.

\paragraph{The v2 pipeline (\fn{generate\_spectral\_data\_v2.jl}).}
\begin{enumerate}
\item $U$ is never formed: $Uv=L_B(L_Av)$ with the two sparse layers (a few MB at $N=14$).
\item Krylov sectors come from a union-find over the supports of the \emph{individual gates}. For $V=H$ this is the same
  partition as the components of $\mathrm{supp}(U)$ (R-T19 Prop.~4), and it is asserted against the old construction
  for every $N\le10$ run. For $V=X$, $U$ is a permutation and is formed directly.
\item Symmetry checks are random-vector tests through the layers. Dense sector blocks are assembled column-chunk-wise,
  and the products inside each block are computed per batch of at most 2048 columns.
\end{enumerate}
\emph{Validation:} v2 reproduces v1 block by block (sector size, block size, sorted quasienergies, $S_2^\mathrm{Haar}$,
degeneracies) for seven rules at $N=9$. Its ring spectra match the exact Python engine to $<10^{-14}$, for both gates.
Non-degenerate eigenvectors agree to $10^{-15}$ in $S_\mathrm{vN}$; degenerate ones are basis dependent by construction.

\paragraph{Analysis cost.} The first analysis pass looped over blocks with one boolean mask per block, which is
$O(\text{blocks}\times\text{states})$: about $10^{11}$ operations for the $k=1$ rules at $N=20$. It now groups the
states by a single sort (\fn{stats.groups}). The full summary of 173 obc0 units takes 70~s, and its output is identical
to the old one on every unit that was not rerun.


\appendix
\section{Bibliography checks}
Every entry below was verified on 2026-10-14 (second pass, Sec.~\ref{sec:lit-known}: 2026-10-15) against the arXiv listing, the DOI or publisher page, or an institutional
repository record. The journal references of Regnault--Nandkishore and Ponte \emph{et al.}\ (PRL 114, 140401) were
confirmed from the arXiv journal-ref field. Deutsch (1991) was confirmed through citing papers only (Rigol \emph{et al.}\
2008 and the reviews). Jordan (1875) is cited for the origin of the two-subspace decomposition; the modern statement is
Halmos (1969).

% R32 bibliography.  Every entry was checked on 2026-10-14 against the arXiv
% listing, the publisher/DOI page or an institutional repository record
% (journal, volume, article/page, year); see R32 appendix "Bibliography checks".
\begin{thebibliography}{99}

\bibitem{atas2013} Y.~Y.~Atas, E.~Bogomolny, O.~Giraud and G.~Roux,
  \emph{Distribution of the ratio of consecutive level spacings in random matrix
  ensembles}, Phys.\ Rev.\ Lett.\ \textbf{110}, 084101 (2013),
  arXiv:1212.5611.

\bibitem{oganesyan2007} V.~Oganesyan and D.~A.~Huse,
  \emph{Localization of interacting fermions at high temperature},
  Phys.\ Rev.\ B \textbf{75}, 155111 (2007), arXiv:cond-mat/0610854.

\bibitem{dyson1962} F.~J.~Dyson,
  \emph{Statistical theory of the energy levels of complex systems.~I},
  J.\ Math.\ Phys.\ \textbf{3}, 140 (1962).

\bibitem{bgs1984} O.~Bohigas, M.~J.~Giannoni and C.~Schmit,
  \emph{Characterization of chaotic quantum spectra and universality of level
  fluctuation laws}, Phys.\ Rev.\ Lett.\ \textbf{52}, 1 (1984).

\bibitem{berrytabor1977} M.~V.~Berry and M.~Tabor,
  \emph{Level clustering in the regular spectrum},
  Proc.\ R.\ Soc.\ Lond.\ A \textbf{356}, 375 (1977).

\bibitem{regnault2016} N.~Regnault and R.~Nandkishore,
  \emph{Floquet thermalization: Symmetries and random matrix ensembles},
  Phys.\ Rev.\ B \textbf{93}, 104203 (2016), arXiv:1510.07653.

\bibitem{giraud2022} O.~Giraud, N.~Mac\'e, \'E.~Vernier and F.~Alet,
  \emph{Probing symmetries of quantum many-body systems through gap ratio
  statistics}, Phys.\ Rev.\ X \textbf{12}, 011006 (2022), arXiv:2008.11173.

\bibitem{dalessio2014} L.~D'Alessio and M.~Rigol,
  \emph{Long-time behavior of isolated periodically driven interacting lattice
  systems}, Phys.\ Rev.\ X \textbf{4}, 041048 (2014), arXiv:1402.5141.

\bibitem{lazarides2014} A.~Lazarides, A.~Das and R.~Moessner,
  \emph{Equilibrium states of generic quantum systems subject to periodic
  driving}, Phys.\ Rev.\ E \textbf{90}, 012110 (2014), arXiv:1403.2946.

\bibitem{ponte2015annphys} P.~Ponte, A.~Chandran, Z.~Papi\'c and D.~A.~Abanin,
  \emph{Periodically driven ergodic and many-body localized quantum systems},
  Ann.\ Phys.\ \textbf{353}, 196 (2015), arXiv:1403.6480.

\bibitem{ponte2015prl} P.~Ponte, Z.~Papi\'c, F.~Huveneers and D.~A.~Abanin,
  \emph{Many-body localization in periodically driven systems},
  Phys.\ Rev.\ Lett.\ \textbf{114}, 140401 (2015), arXiv:1410.8518.

\bibitem{deutsch1991} J.~M.~Deutsch,
  \emph{Quantum statistical mechanics in a closed system},
  Phys.\ Rev.\ A \textbf{43}, 2046 (1991).

\bibitem{srednicki1994} M.~Srednicki, \emph{Chaos and quantum thermalization},
  Phys.\ Rev.\ E \textbf{50}, 888 (1994), arXiv:cond-mat/9403051.

\bibitem{rigol2008} M.~Rigol, V.~Dunjko and M.~Olshanii,
  \emph{Thermalization and its mechanism for generic isolated quantum systems},
  Nature \textbf{452}, 854 (2008), arXiv:0708.1324.

\bibitem{dalessio2016} L.~D'Alessio, Y.~Kafri, A.~Polkovnikov and M.~Rigol,
  \emph{From quantum chaos and eigenstate thermalization to statistical
  mechanics and thermodynamics}, Adv.\ Phys.\ \textbf{65}, 239 (2016),
  arXiv:1509.06411.

\bibitem{page1993} D.~N.~Page, \emph{Average entropy of a subsystem},
  Phys.\ Rev.\ Lett.\ \textbf{71}, 1291 (1993), arXiv:gr-qc/9305007.

\bibitem{vidmar2017} L.~Vidmar and M.~Rigol,
  \emph{Entanglement entropy of eigenstates of quantum chaotic Hamiltonians},
  Phys.\ Rev.\ Lett.\ \textbf{119}, 220603 (2017), arXiv:1708.08453.

\bibitem{bernien2017} H.~Bernien, S.~Schwartz, A.~Keesling, H.~Levine,
  A.~Omran, H.~Pichler, S.~Choi, A.~S.~Zibrov, M.~Endres, M.~Greiner,
  V.~Vuleti\'c and M.~D.~Lukin,
  \emph{Probing many-body dynamics on a 51-atom quantum simulator},
  Nature \textbf{551}, 579 (2017), arXiv:1707.04344.

\bibitem{turner2018} C.~J.~Turner, A.~A.~Michailidis, D.~A.~Abanin,
  M.~Serbyn and Z.~Papi\'c,
  \emph{Weak ergodicity breaking from quantum many-body scars},
  Nat.\ Phys.\ \textbf{14}, 745 (2018).

\bibitem{shiraishi2017} N.~Shiraishi and T.~Mori,
  \emph{Systematic construction of counterexamples to the eigenstate
  thermalization hypothesis}, Phys.\ Rev.\ Lett.\ \textbf{119}, 030601 (2017),
  arXiv:1702.08227.

\bibitem{mondaini2018} R.~Mondaini, K.~Mallayya, L.~F.~Santos and M.~Rigol,
  \emph{Comment on ``Systematic construction of counterexamples to the
  eigenstate thermalization hypothesis''}, Phys.\ Rev.\ Lett.\ \textbf{121},
  038901 (2018), arXiv:1711.06279.

\bibitem{serbyn2021} M.~Serbyn, D.~A.~Abanin and Z.~Papi\'c,
  \emph{Quantum many-body scars and weak breaking of ergodicity},
  Nat.\ Phys.\ \textbf{17}, 675 (2021), arXiv:2011.09486.

\bibitem{moudgalya2022review} S.~Moudgalya, B.~A.~Bernevig and N.~Regnault,
  \emph{Quantum many-body scars and Hilbert space fragmentation: a review of
  exact results}, Rep.\ Prog.\ Phys.\ \textbf{85}, 086501 (2022),
  arXiv:2109.00548.

\bibitem{chandran2023} A.~Chandran, T.~Iadecola, V.~Khemani and R.~Moessner,
  \emph{Quantum many-body scars: A quasiparticle perspective},
  Annu.\ Rev.\ Condens.\ Matter Phys.\ \textbf{14}, 443 (2023),
  arXiv:2206.11528.

\bibitem{sugiura2021} S.~Sugiura, T.~Kuwahara and K.~Saito,
  \emph{Many-body scar state intrinsic to periodically driven system},
  Phys.\ Rev.\ Research \textbf{3}, L012010 (2021), arXiv:1911.06092.

\bibitem{sala2020} P.~Sala, T.~Rakovszky, R.~Verresen, M.~Knap and
  F.~Pollmann, \emph{Ergodicity breaking arising from Hilbert space
  fragmentation in dipole-conserving Hamiltonians}, Phys.\ Rev.\ X
  \textbf{10}, 011047 (2020), arXiv:1904.04266.

\bibitem{khemani2020} V.~Khemani, M.~Hermele and R.~Nandkishore,
  \emph{Localization from Hilbert space shattering: From theory to physical
  realizations}, Phys.\ Rev.\ B \textbf{101}, 174204 (2020), arXiv:1910.01137.

\bibitem{moudgalya2022} S.~Moudgalya and O.~I.~Motrunich,
  \emph{Hilbert space fragmentation and commutant algebras},
  Phys.\ Rev.\ X \textbf{12}, 011050 (2022), arXiv:2108.10324.

\bibitem{kos2018} P.~Kos, M.~Ljubotina and T.~Prosen,
  \emph{Many-body quantum chaos: Analytic connection to random matrix theory},
  Phys.\ Rev.\ X \textbf{8}, 021062 (2018), arXiv:1712.02665.

\bibitem{chan2018} A.~Chan, A.~De~Luca and J.~T.~Chalker,
  \emph{Solution of a minimal model for many-body quantum chaos},
  Phys.\ Rev.\ X \textbf{8}, 041019 (2018), arXiv:1712.06836.

\bibitem{gopalakrishnan2018} S.~Gopalakrishnan and B.~Zakirov,
  \emph{Facilitated quantum cellular automata as simple models with non-thermal
  eigenstates and dynamics}, Quantum Sci.\ Technol.\ \textbf{3}, 044004 (2018),
  arXiv:1802.07729.

\bibitem{iadecola2020} T.~Iadecola and S.~Vijay,
  \emph{Nonergodic quantum dynamics from deformations of classical cellular
  automata}, Phys.\ Rev.\ B \textbf{102}, 180302 (2020), arXiv:2006.02440.

\bibitem{buca2021} B.~Bu\v{c}a, K.~Klobas and T.~Prosen,
  \emph{Rule 54: exactly solvable model of nonequilibrium statistical
  mechanics}, J.\ Stat.\ Mech.\ (2021) 074001, arXiv:2103.16543.

\bibitem{wilkinson2020} J.~W.~P.~Wilkinson, K.~Klobas, T.~Prosen and
  J.~P.~Garrahan, \emph{Exact solution of the Floquet-PXP cellular automaton},
  Phys.\ Rev.\ E \textbf{102}, 062107 (2020), arXiv:2006.06556.

\bibitem{jordan1875} C.~Jordan, \emph{Essai sur la g\'eom\'etrie \`a $n$
  dimensions}, Bull.\ Soc.\ Math.\ France \textbf{3}, 103 (1875).

\bibitem{halmos1969} P.~R.~Halmos, \emph{Two subspaces},
  Trans.\ Amer.\ Math.\ Soc.\ \textbf{144}, 381 (1969).

\end{thebibliography}

\end{document}
